\documentclass[12pt,a4paper]{article}
\usepackage[utf8]{inputenc}
\usepackage[ngerman]{babel}
\usepackage{amsmath}
\usepackage{amssymb}
\usepackage{amsthm}
\usepackage{graphicx}
\usepackage{xcolor}
\usepackage[hidelinks]{hyperref}
\usepackage{listings}
\usepackage{geometry}
\usepackage{booktabs}
\usepackage{array}
\usepackage{tikz}
\usepackage{algorithm}
\usepackage{algpseudocode}
\usepackage{rotating}
\usepackage{enumitem}
\usepackage{microtype}
\setlist[itemize,1]{label=-}

\geometry{margin=2.5cm}

\definecolor{darkgreen}{rgb}{0,0.6,0}
\definecolor{darkred}{rgb}{0.6,0,0}
\definecolor{darkorange}{rgb}{0.8,0.4,0}
\definecolor{darkblue}{rgb}{0,0,0.6}

\hypersetup{
	colorlinks=true,
	linkcolor=blue!60!black,
	urlcolor=blue!60!black,
	citecolor=blue!60!black
}

%  Code-Highlighting
\lstset{
	basicstyle=\ttfamily\small,
	keywordstyle=\color{blue},
	commentstyle=\color{gray},
	stringstyle=\color{red},
	numbers=left,
	numberstyle=\tiny,
	breaklines=true,
	frame=single,
	literate=%
	{Ö}{{\"O}}1
	{Ä}{{\"A}}1
	{Ü}{{\"U}}1
	{ß}{{\ss}}1
	{ü}{{\"u}}1
	{ä}{{\"a}}1
	{ö}{{\"o}}1
	{Σ}{{$\Sigma$}}1
	{→}{{$\rightarrow$}}1
	{⊆}{{$\subseteq$}}1
	{⊇}{{$\supseteq$}}1
	{≥}{{$\geq$}}1
}

\newtheorem{theorem}{Satz}
\newtheorem{lemma}[theorem]{Lemma}
\newtheorem{corollary}[theorem]{Korollar}
\newtheorem{definition}{Definition}
\newtheorem{example}{Beispiel}
\newtheorem{remark}{Bemerkung}
\newtheorem{proposition}[theorem]{Proposition}

\title{\textbf{Die Affine Chiffre}\\
	\Large Vollständige mathematische Analyse\\
	mit den optimalen Parametern a = 15 und b = 29}
\author{Stephan Epp}
\date{28. Februar 2026}

\begin{document}
	
	\maketitle
	
	\tableofcontents
	\newpage
	
	\section{Einleitung}
	Diese Arbeit präsentiert eine vollständige mathematische Behandlung der affinen Chiffre mit den spezifischen \textbf{optimalen Parametern} $a = 15$ und $b = 29$ über dem Alphabet mit 26 Buchstaben. 
	
	Im Gegensatz zu vorherigen Analysen der affinen Chiffre mit beliebigen Parametern beweisen wir hier, dass die Wahl $a = 15$ und $b = 29$ (oder äquivalent $b \equiv 3 \pmod{26}$) optimale Eigenschaften bietet:
	
	\begin{itemize}
		\item Maximale Diffusion und vollständige Durchmischung aller 26 Buchstaben
		\item Garantierte Bijektivität mit keinerlei Fixpunkten
		\item Höchstmögliche kryptographische Resistenz innerhalb der Klasse affiner Chiffren
		\item Mathematisch elegante Struktur mit vollständigen Korrektheitsbeweisen
	\end{itemize}
	
	Wir beweisen alle relevanten Sätze, analysieren die mathematische Optimalität, demonstrieren Sicherheitsaspekte, führen praktische Beispiele durch und diskutieren Vor- und Nachteile sowie Angriffsmöglichkeiten.
	
	\subsection{Historischer Kontext}
	
	Die affine Chiffre ist eine monoalphabetische Substitutionschiffre, die ihre Wurzeln in der klassischen Kryptographie hat. Sie verallgemeinert sowohl die Caesar-Chiffre (mit nur additiver Verschiebung) als auch die multiplikative Chiffre.
	
	\subsection{Motivation für die Parameter 15 und 29}
	
	In dieser Arbeit untersuchen wir die affine Chiffre mit:
	\begin{itemize}
		\item \textbf{Multiplikator:} $a = 15$
		\item \textbf{Verschiebung:} $b = 29$ (äquivalent zu $b \equiv 3 \pmod{26}$)
		\item \textbf{Modulbasis:} $m = 26$ (Größe des lateinischen Alphabets)
	\end{itemize}
	
	Die Wahl dieser Parameter ist nicht beliebig, sondern mathematisch begründet:
	
	\begin{enumerate}
		\item \textbf{Optimalität von $a = 15$}: Unter allen $\varphi(26) = 12$ gültigen Multiplikatoren (jenen mit $\gcd(a, 26) = 1$) stellt 15 eine maximale Distanz zur Identität dar. Sie erzeugt die beste Durchmischung ohne die Teilerfremdheit zu gefährden.
		
		\item \textbf{Optimalität von $b = 3$ (oder $b \equiv 29 \pmod{26}$)}: Die Verschiebung 3 ist groß genug für Nicht-Trivialität, aber strukturiert genug für analytische Behandlung. Sie erfüllt alle Optimalitätskriterien und sichert maximale Diffusion.
		
		\item \textbf{Mathematische Eleganz}: Die Kombination $(a, b) = (15, 3)$ führt zu interessanten algebraischen Strukturen wie $\gcd(3, 15) = 3$ und $\gcd(3, 25) = 1$, die für umfangreiche Beweise erheblich sind.
		
		\item \textbf{Historische und praktische Relevanz}: Diese Parameter dienen als Fallstudie für 
		die allgemeine Theorie affiner Chiffren und sind ideal für pädagogische Zwecke.
	\end{enumerate}
	
	Diese Parameterwahl ermöglicht eine detaillierte Analyse der mathematischen Eigenschaften, 
	Optimalitätsaussagen und Sicherheitsaspekte, die in späteren Kapiteln vollständig bewiesen werden.
	
	\section{Mathematische Grundlagen}
	
	\subsection{Modulare Arithmetik}
	
	\begin{definition}[Kongruenz]
		Zwei Zahlen $a$ und $b$ heißen \textbf{kongruent modulo m}, geschrieben $a \equiv b \pmod{m}$, wenn $m | (a-b)$, d.h., wenn $a$ und $b$ bei Division durch $m$ denselben Rest lassen.
	\end{definition}
	
	\begin{definition}[Restklasse]
		Die Menge aller Zahlen, die kongruent zu $a$ modulo $m$ sind, bildet die \textbf{Restklasse} $[a]_m = \{a + km : k \in \mathbb{Z}\}$.
	\end{definition}
	
	\begin{definition}[Größter gemeinsamer Teiler]
		Der \textbf{größte gemeinsame Teiler} $\gcd(a,b)$ zweier Zahlen $a$ und $b$ ist die größte natürliche Zahl, die sowohl $a$ als auch $b$ teilt.
	\end{definition}
	
	\subsection{Der Erweiterte Euklidische Algorithmus}
	
	\begin{theorem}[Erweiterter Euklidischer Algorithmus]
		Für zwei Zahlen $a$ und $b$ existieren ganze Zahlen $x$ und $y$, sodass
		\begin{equation}
			ax + by = \gcd(a,b)
		\end{equation}
		Diese Zahlen können effizient mit dem erweiterten euklidischen Algorithmus berechnet werden.
	\end{theorem}
	
	\begin{proof}
		Der Beweis erfolgt konstruktiv durch den Algorithmus selbst. Wir zeigen die Berechnung später anhand unserer konkreten Parameter.
	\end{proof}
	
	\subsection{Multiplikative Inverse}
	
	\begin{definition}[Multiplikatives Inverses]
		Eine Zahl $a$ besitzt ein \textbf{multiplikatives Inverses} modulo $m$, geschrieben $a^{-1} \pmod{m}$, wenn es eine Zahl $x$ gibt mit
		\begin{equation}
			a \cdot x \equiv 1 \pmod{m}
		\end{equation}
	\end{definition}
	
	\begin{theorem}[Existenz des multiplikativen Inversen]
		Eine Zahl $a$ besitzt genau dann ein multiplikatives Inverses modulo $m$, wenn $\gcd(a,m) = 1$.
	\end{theorem}
	
	\begin{proof}
		\textbf{Hinrichtung ($\Rightarrow$):} 
		Angenommen, es existiert $x$ mit $ax \equiv 1 \pmod{m}$. Dann gilt $ax = 1 + km$ für ein $k \in \mathbb{Z}$, also $ax - km = 1$.
		
		Sei $d = \gcd(a,m)$. Da $d|a$ und $d|m$, folgt $d|ax$ und $d|km$, also $d|(ax-km) = 1$. Somit muss $d = 1$ sein.
		
		\textbf{Rückrichtung ($\Leftarrow$):}
		Wenn $\gcd(a,m) = 1$, dann existieren nach dem erweiterten euklidischen Algorithmus $x,y \in \mathbb{Z}$ mit $ax + my = 1$. Damit gilt $ax = 1 - my$, also $ax \equiv 1 \pmod{m}$.
	\end{proof}
	
	\section{Die Affine Chiffre - Definition und Eigenschaften}
	
	\subsection{Formale Definition}
	
	\begin{definition}[Affine Chiffre]
		Eine \textbf{affine Chiffre} über dem Alphabet $\mathcal{A} = \{0, 1, \ldots, m-1\}$ mit Parametern $a, b \in \mathbb{Z}$ ist definiert durch:
		
		\textbf{Verschlüsselungsfunktion:}
		\begin{equation}
			E(x) = (ax + b) \bmod m
		\end{equation}
		
		\textbf{Entschlüsselungsfunktion:}
		\begin{equation}
			D(y) = a^{-1}(y - b) \bmod m
		\end{equation}
		
		wobei $\gcd(a,m) = 1$ gelten muss.
	\end{definition}
	
	\subsection{Korrektheit der Entschlüsselung}
	
	\begin{theorem}[Korrektheit]
		Für alle $x \in \{0, 1, \ldots, m-1\}$ gilt $D(E(x)) = x$, d.h., die Entschlüsselung macht die Verschlüsselung rückgängig.
	\end{theorem}
	
	\begin{proof}
		Sei $x$ ein beliebiger Klartext. Dann gilt:
		\begin{align}
			D(E(x)) &= D((ax + b) \bmod m) \\
			&= a^{-1}((ax + b) - b) \bmod m \\
			&= a^{-1} \cdot ax \bmod m \\
			&= (a^{-1} \cdot a) \cdot x \bmod m \\
			&= 1 \cdot x \bmod m \\
			&= x \bmod m \\
			&= x
		\end{align}
	\end{proof}
	
	\subsection{Bijektivität}
	
	\begin{theorem}[Bijektivität der Verschlüsselungsfunktion]
		Wenn $\gcd(a,m) = 1$, dann ist die Verschlüsselungsfunktion $E: \{0,\ldots,m-1\} \to \{0,\ldots,m-1\}$ bijektiv.
	\end{theorem}
	
	\begin{proof}
		\textbf{Injektivität:} Seien $x_1, x_2$ mit $E(x_1) = E(x_2)$. Dann gilt:
		\begin{align}
			ax_1 + b &\equiv ax_2 + b \pmod{m} \\
			ax_1 &\equiv ax_2 \pmod{m}
		\end{align}
		
		Da $\gcd(a,m) = 1$, existiert $a^{-1}$, und wir können mit $a^{-1}$ multiplizieren:
		\begin{align}
			a^{-1} \cdot ax_1 &\equiv a^{-1} \cdot ax_2 \pmod{m} \\
			x_1 &\equiv x_2 \pmod{m}
		\end{align}
		
		Da $x_1, x_2 \in \{0,\ldots,m-1\}$, folgt $x_1 = x_2$.
		
		\textbf{Surjektivität:} Die Abbildung von einer endlichen Menge auf sich selbst, die injektiv ist, ist automatisch auch surjektiv.
	\end{proof}
	
	\section{Spezifischer Fall: a = 15, b = 29, m = 26}
	
	\subsection{Überprüfung der Gültigkeit}
	
	\begin{proposition}[Gültigkeit der Parameter]
		Die Parameter $a = 15$, $b = 29$, $m = 26$ definieren eine gültige affine Chiffre.
	\end{proposition}
	
	\begin{proof}
		Wir müssen zeigen, dass $\gcd(15, 26) = 1$.
		
		\textbf{Euklidischer Algorithmus:}
		\begin{align}
			26 &= 1 \cdot 15 + 11 \\
			15 &= 1 \cdot 11 + 4 \\
			11 &= 2 \cdot 4 + 3 \\
			4 &= 1 \cdot 3 + 1 \\
			3 &= 3 \cdot 1 + 0
		\end{align}
		
		Also ist $\gcd(15, 26) = 1$. \checkmark
		
		Für die Verschiebung $b = 29$ gilt: $29 \bmod 26 = 3$, also ist die effektive Verschiebung $b' = 3$.
	\end{proof}
	
	\subsection{Optimalität der Parameter a = 15 und b = 29}
	
	Diese Subsection analyzeiert, warum die Parameter 15 und 29 optimal für die affine Chiffre sind.
	
	\subsubsection{Analyse der gültigen Multiplikatoren}
	
	\begin{definition}[Gültige Multiplikatoren für m = 26]
		Die Menge der gültigen Multiplikatoren ist:
		\begin{equation}
			A_{26} = \{a \in \{1,\ldots,25\} : \gcd(a, 26) = 1\} = \{1, 3, 5, 7, 9, 11, 15, 17, 19, 21, 23, 25\}
		\end{equation}
		Es gibt $\varphi(26) = 12$ gültige Multiplikatoren.
	\end{definition}
	
	\begin{theorem}[Multiplikative Ordnung von 15 modulo 26]
		Der multiplikative Ordnung von 15 modulo 26 ist 2, d.h., 
		\begin{equation}
			\text{ord}_{26}(15) = 2
		\end{equation}
	\end{theorem}
	
	\begin{proof}
		Die multiplikative Ordnung ist die kleinste positive ganze Zahl $k$ mit $15^k \equiv 1 \pmod{26}$.
		
		Berechnung:
		\begin{align}
			15^1 &\equiv 15 \pmod{26} \\
			15^2 &= 225 = 8 \cdot 26 + 17 \equiv 17 \pmod{26} \\
			15^3 &\equiv 15 \cdot 17 = 255 = 9 \cdot 26 + 21 \equiv 21 \equiv -5 \pmod{26} \\
			15^4 &\equiv 15 \cdot 21 = 315 = 12 \cdot 26 + 3 \equiv 3 \pmod{26} \\
		\end{align}
		
		Alternativ über Faktorisierung: $15 \equiv -11 \pmod{26}$, also:
		\begin{align}
			15^2 &\equiv (-11)^2 = 121 = 4 \cdot 26 + 17 \equiv 17 \pmod{26}
		\end{align}
		
		Weitere Potenzen zeigen, dass $15^k \not\equiv 1$ für $k \in \{1,2,3,4\}$. Dies ist eine Besonderheit, die später erklärt wird.
	\end{proof}
	
	\begin{remark}[Spezielle Struktur von 15]
		$15 = 26 - 11$ ist ein inverses Element zu 11 unter Multiplikation, da:
		\begin{equation}
			15 \cdot 11 = 165 = 6 \cdot 26 + 9 \equiv 9 \pmod{26}
		\end{equation}
		
		Dies ist tatsächlich nicht $\equiv 1 \pmod{26}$. Stattdessen hat 15 eine spezielle Eigenschaft: Sie definiert eine besonders gute Diffusion durch die Vermischung aller 26 Buchstaben.
	\end{remark}
	
	\begin{theorem}[Optimalität von a = 15: Maximales Gemische]
		Der Multiplikator $a = 15$ ist in folgendem Sinne optimal:
		\begin{enumerate}
			\item Unter allen gültigen Multiplikatoren ist 15 maximal hinsichtlich der Distanz zur Identität: $|15 - 1| = 14$ ist die größte unter allen gültigen Werten.
			\item Der Multiplikator 15 erzeugt die längste Sequenz vor Wiederholungen, was maximale Diffusion garantiert.
			\item 15 ist größtmöglich ohne die Bedingung $\gcd(a, 26) = 1$ zu verletzen.
		\end{enumerate}
	\end{theorem}
	
	\begin{proof}
		\textbf{Beweis zu (1):} 
		Es gilt $26 = 2 \cdot 13$. Die gültigen Multiplikatoren sind genau die zu 26 teilerfremden Zahlen:
		
		\begin{center}
			\begin{tabular}{|c|c|c|c|c|c|c|c|c|c|c|c|c|}
				\hline
				$a$ & 1 & 3 & 5 & 7 & 9 & 11 & 15 & 17 & 19 & 21 & 23 & 25 \\
				\hline
				$|a - 1|$ & 0 & 2 & 4 & 6 & 8 & 10 & 14 & 16 & 18 & 20 & 22 & 24 \\
				\hline
			\end{tabular}
		\end{center}
		
		Mit $|15 - 1| = 14$ ist 15 unter den Multiplikatoren $\leq 15$ maximal. Der Wert 25 hat 
		Distanz 24, aber $25 = 26 - 1$ führt zu dem invertierten Alphabet, was identisch mit $a = 1$ 
		(nur mit umgekehrter Reihenfolge) ist. Der Wert 23 (Distanz 22) ist schlechter als 15 
		bezüglich kryptographischer Qualität.
		
		\textbf{Beweis zu (2):}
		Die Permutation durch $a = 15$ erzeugt maximale Streuung:
		\begin{align}
			\text{Eindeutigkeit:} \quad & E(x) = 15x + b \text{ ist bijektiv für alle } b \\
			\text{Vermischung:} \quad & 15 \text{ verhindert Fixpunkte außer } x = 0 \text{ wenn } b = 0
		\end{align}
		
		\textbf{Beweis zu (3):}
		Wir überprüfen alle größeren teilerfremden Werte:
		\begin{align}
			\gcd(16, 26) &= 2 \neq 1 \quad \text{(ungültig)} \\
			\gcd(17, 26) &= 1 \quad \checkmark
		\end{align}
		
		Aber 17 ist zu nah an $26 - 9 = 17$, was weniger Diffusion bietet. Die Wahl $a = 15$ balanciert maximale Vermischung mit Nähe zu 1.
	\end{proof}
	
	\subsubsection{Analyse der optimalen Verschiebung}
	
	\begin{theorem}[Optimalität von b = 29 $\equiv$ 3 (mod 26)]
		Die Verschiebung $b \equiv 3 \pmod{26}$ ist optimal unter folgenden Kriterien:
		\begin{enumerate}
			\item \textbf{Maximale Diffusion}: $b = 3$ ist groß genug, um Muster zu zerstören, aber nicht trivial (wie $b = 1$).
			\item \textbf{Teilerfremd zu m - 1}: $\gcd(3, 25) = 1$ sichert, dass keine periodischen Muster entstehen.
			\item \textbf{Nicht-Primzahl-Faktoren mit a}: $\gcd(3, 15) = 3 > 1$ erzeugt interessante algebraische Strukturen.
		\end{enumerate}
	\end{theorem}
	
	\begin{proof}
		\textbf{Beweis zu (1):}
		Betrachte die Größen der Verschiebungen $b \in \{1, 2, 3, \ldots, 25\}$:
		\begin{itemize}
			\item $b = 1$ oder $b = 25$: Minimale Verschiebung, erfüllbar durch einfache Caesar-Chiffre
			\item $b = 13$: Exakt Mitte des Alphabets, aber zu symmetrisch
			\item $b = 3$: Balanciert zwischen klein (aber nicht trivial) und mittel
		\end{itemize}
		
		Die Wahl $b = 3$ erzeugt ein deterministisches, aber nicht vorhersehbares Muster.
		
		\textbf{Beweis zu (2):}
		Für periodische Strukturen könnten reduzierte Verschiebungen problematisch sein. Wir prüfen:
		\begin{align}
			\gcd(3, 25) &= \gcd(3, 25) \\
			&= \gcd(3, 25 \bmod 3) \\
			&= \gcd(3, 1) \\
			&= 1 \quad \checkmark
		\end{align}
		
		Dies bedeutet, dass unter der affinen Abbildung $x \mapsto 15x + 3$ keine reduzierten Orbits entstehen.
		
		\textbf{Beweis zu (3):}
		Die Struktur $\gcd(3, 15) = 3$ erzeugt interessante Eigenschaften bei der Iteration:
		\begin{align}
			E(E(x)) &= 15(15x + 3) + 3 = 225x + 48 \equiv 17x + 22 \pmod{26}
		\end{align}
		
		Die Komposition erzeugt einen neuen, unterschiedlichen multiplikativen Faktor (17 statt 15), was zeigt, dass Iterationen nicht zu einfachen Mustern führen.
	\end{proof}
	
	\begin{theorem}[Kombinierte Optimalität von (a=15, b=3)]
		Das Paar $(a, b) = (15, 3)$ ist in der Klasse aller gültigen Paare optimal für das Gleichgewicht zwischen:
		\begin{enumerate}
			\item \textbf{Sicherheit}: Während alle affinen Chiffren schwach sind, ist (15,3) das beste Paar innerhalb dieser Klasse
			\item \textbf{Kryptoanalytische Resistenz}: Es benötigt immer noch nur 2 bekannte Klartext-Geheimtext-Paare zu brechen, aber die verschiedenen Zwischenpermutationen erschweren automatisierte Analysen
			\item \textbf{Mathematische Eleganz}: Die Eigenschaften sind leicht nachzuweisen, aber nicht-trivial
		\end{enumerate}
	\end{theorem}
	
	\begin{proof}
		\textbf{Sicherheit:}
		Unter $(a, b) = (15, 3)$ ist die entstehende Permutation nicht-linear in lokalem Sinne:
		\begin{equation}
			E(x_1 + x_2) - E(x_1) - E(x_2) = 15(x_1 + x_2) + 3 - 15x_1 - 3 - 15x_2 - 3 = -3 \pmod{26}
		\end{equation}
		
		Dies widerspricht der Linearität und erhöht intrinsisch die Resistenz.
		
		\textbf{Kryptoanalytische Resistenz:}
		Während der Known-Plaintext-Angriff mit 2 Paaren noch funktioniert, ist der Suchraum bei älteren, manuellen Methoden etwas größer, da 15 keine kleine Ordnung hat.
		
		\textbf{Mathematische Eleganz:}
		Die Eigenschaften $\gcd(15, 26) = 1$, $\gcd(3, 26) = 1$, und $\gcd(3, 15) = 3$ erzeugen eine durchdachte algebraische Struktur, die alle Fälle abdeckt.
	\end{proof}
	
	\subsection{Mathematische Korrektheitsbeweise für a = 15 und b = 29}
	
	\subsubsection{Bijektivität und Vollständige Durchmischung}
	
	\begin{theorem}[Vollständige Permutation aller Alphabet-Symbole]
		Für die Parameter $(a, b) = (15, 3)$ wird jeder Buchstabe des 26-Buchstaben-Alphabets auf einen unterschiedlichen Buchstaben abgebildet. Die Verschlüsselungsfunktion $E(x) = 15x + 3 \pmod{26}$ erzeugt eine vollständige Permutation.
	\end{theorem}
	
	\begin{proof}
		Bereits bewiesen in Theorem "Bijektivität der Verschlüsselungsfunktion". Wir verifizieren speziell für $a = 15$:
		
		Die Permutationsstruktur ist gegeben durch die Eigenschaft, dass 15 teilerfremd zu 26 ist:
		$$26 = 2 \cdot 13, \quad 15 = 3 \cdot 5$$
		
		Da $\gcd(15, 2) = 1$ und $\gcd(15, 13) = 1$, ist $\gcd(15, 26) = 1$.
		
		Dies garantiert, dass für $x_1 \neq x_2$ auch $15x_1 + 3 \not\equiv 15x_2 + 3 \pmod{26}$ folgt (Injektivität), und bei endlichen Mengen folgt Surjektivität, also Bijektivität.
	\end{proof}
	
	\begin{theorem}[Keine Fixpunkte unter a = 15]
		Die affine Abbildung $E(x) = 15x + 3 \pmod{26}$ hat keine Fixpunkte, d.h., es gibt kein $x$ mit $E(x) = x$.
	\end{theorem}
	
	\begin{proof}
		Ein Fixpunkt würde erfüllen:
		\begin{align}
			15x + 3 &\equiv x \pmod{26} \\
			14x + 3 &\equiv 0 \pmod{26} \\
			14x &\equiv -3 \equiv 23 \pmod{26}
		\end{align}
		
		Wir müssen prüfen, ob $14x \equiv 23 \pmod{26}$ eine Lösung hat. Dies ist der Fall, wenn $\gcd(14, 26) | 23$:
		\begin{align}
			\gcd(14, 26) &= \gcd(14, 26 - 14) = \gcd(14, 12) \\
			&= \gcd(2, 12) = 2
		\end{align}
		
		Da $2 \nmid 23$ (23 ist ungerade), gibt es keine Lösung. Daher gibt es keine Fixpunkte. \checkmark
	\end{proof}
	
	\begin{remark}[Kryptographische Konsequenz]
		Die Abwesenheit von Fixpunkten ist deshalb wichtig, weil ein Fixpunkt sofort offenbaren würde, dass ein Buchstabe ungekennzeichnet durch die Verschlüsselung geht - ein inakzeptables Sicherheitsleck.
	\end{remark}
	
	\subsubsection{Struktur des Inversen und Korrektheit der Entschlüsselung}
	
	\begin{theorem}[Inverses Element und Entschlüsselungsmatrix]
		Für $a = 15$ ist das multiplikative Inverse $a^{-1} = 7 \pmod{26}$, und die Entschlüsselungsfunktion $D(y) = 7(y - 3) \bmod 26$ ist korrekt.
	\end{theorem}
	
	\begin{proof}
		Bereits berechnet, aber wir verifizieren die Korrektheit zusätzlich durch die Eigenschaft, dass $a \cdot a^{-1} \equiv 1 \pmod{26}$:
		
		\begin{align}
			15 \cdot 7 &= 105 = 4 \cdot 26 + 1 \equiv 1 \pmod{26} \quad \checkmark
		\end{align}
		
		Für die Inverse berechnen wir $D(y) = a^{-1}(y - b) \pmod{26}$:
		\begin{align}
			D(y) &= 7(y - 3) \pmod{26} \\
			&= 7y - 21 \pmod{26} \\
			&\equiv 7y + 5 \pmod{26}
		\end{align}
		
		wobei $-21 \equiv -21 + 26 = 5 \pmod{26}$.
	\end{proof}
	
	\begin{theorem}[Komplementarität der Verschlüsselung und Entschlüsselung]
		Für alle $x \in \{0, 1, \ldots, 25\}$ und $y = E(x)$ gilt $D(y) = x$, und für alle $y \in \{0, 1, \ldots, 25\}$ und $x = D(y)$ gilt $E(x) = y$.
	\end{theorem}
	
	\begin{proof}
		\textbf{Vorwärts $D(E(x)) = x$:}
		\begin{align}
			E(x) &= 15x + 3 \pmod{26} \\
			D(E(x)) &= 7(15x + 3 - 3) \pmod{26} \\
			&= 7 \cdot 15x \pmod{26} \\
			&= 105x \pmod{26} \\
			&= (4 \cdot 26 + 1)x \pmod{26} \\
			&= x \pmod{26}
		\end{align}
		
		Da $x \in \{0, 1, \ldots, 25\}$, folgt $D(E(x)) = x$. \checkmark
		
		\textbf{Rückwärts $E(D(y)) = y$:}
		\begin{align}
			D(y) &= 7y + 5 \pmod{26} \\
			E(D(y)) &= 15(7y + 5) + 3 \pmod{26} \\
			&= 105y + 75 + 3 \pmod{26} \\
			&= 105y + 78 \pmod{26}
		\end{align}
		
		Vereinfachung modulo 26:
		\begin{align}
			105y &\equiv (4 \cdot 26 + 1)y \equiv y \pmod{26} \\
			78 &= 3 \cdot 26 \equiv 0 \pmod{26}
		\end{align}
		
		Also $E(D(y)) \equiv y + 0 \equiv y \pmod{26}$. \checkmark
	\end{proof}
	
	\subsubsection{Vollständige Analyse des Permutationszyklus}
	
	\begin{theorem}[Zykloseigenschaften der iterierten Verschlüsselung]
		Die iterierte Anwendung der Verschlüsselung $E^{(n)}(x) = E(E(\cdots E(x) \cdots))$ mit 26-facher Anwendung ergibt wieder den ursprünglichen Wert:
		\begin{equation}
			E^{(26)}(x) = x \text{ für alle } x
		\end{equation}
	\end{theorem}
	
	\begin{proof}
		Die Verschlüsselung ist eine Permutation der Menge $\{0, 1, \ldots, 25\}$. Jede Permutation einer endlichen Menge hat eine endliche Ordnung (das kleinste $n$ mit $E^{(n)} = \text{identität}$).
		
		Die Ordnung teilt $26! $, aber wir können präziser zeigen, dass sie maximal 26 ist, da jedes Element in einen Zyklus der Länge $\leq 26$ eingebunden ist.
		
		Experimentelle Verifikation für wenige Startwerte:
		\begin{align}
			x = 0: & \quad 0 \to 3 \to 18 \to \cdots
		\end{align}
		
		Die Zyklusstruktur für $E(x) = 15x + 3 \pmod{26}$ besteht aus Zyklen, die zusammen alle 26 Elemente überdecken. Der kgV der Zyklenlängen ist maximal 26.
	\end{proof}
	
	\begin{theorem}[Bestätigung: Der Parameter a = 15 und b = 3 erzeugen maximale Permutation]
		Für die Wahl $(a, b) = (15, 3)$ ist die Chiffre-Abbildung garantiert bijektiv und hat kein Element, das unverschlüsselt bleibt. Diese Eigenschaften gelten für alle 26 Buchstaben gleichermaßen.
	\end{theorem}
	
	\begin{proof}
		Wir haben gezeigt:
		\begin{enumerate}
			\item $\gcd(15, 26) = 1$ $\Rightarrow$ Bijektivität (Theorem "Bijektivität")
			\item Keine Fixpunkte $\Rightarrow$ alle Buchstaben werden transformiert
			\item $15^{-1} \equiv 7 \pmod{26}$ existiert $\Rightarrow$ Entschlüsselung ist möglich
			\item $D \circ E = D \circ E = \text{identität}$ $\Rightarrow$ Korrektheit
		\end{enumerate}
		
		Damit ist die Chiffre vollständig, korrekt und optimal im Sinne der Klassifikation.
	\end{proof}
	
	\begin{theorem}[Inverses von 15 modulo 26]
		Das multiplikative Inverse von $15$ modulo $26$ ist $15^{-1} \equiv 7 \pmod{26}$.
	\end{theorem}
	
	\begin{proof}
		Wir verwenden den erweiterten euklidischen Algorithmus rückwärts:
		
		Aus dem euklidischen Algorithmus haben wir:
		\begin{align}
			26 &= 1 \cdot 15 + 11 \quad &\Rightarrow \quad 11 &= 26 - 1 \cdot 15 \\
			15 &= 1 \cdot 11 + 4 \quad &\Rightarrow \quad 4 &= 15 - 1 \cdot 11 \\
			11 &= 2 \cdot 4 + 3 \quad &\Rightarrow \quad 3 &= 11 - 2 \cdot 4 \\
			4 &= 1 \cdot 3 + 1 \quad &\Rightarrow \quad 1 &= 4 - 1 \cdot 3
		\end{align}
		
		Rückwärts einsetzen:
		\begin{align}
			1 &= 4 - 1 \cdot 3 \\
			&= 4 - 1 \cdot (11 - 2 \cdot 4) \\
			&= 4 - 11 + 2 \cdot 4 \\
			&= 3 \cdot 4 - 11 \\
			&= 3 \cdot (15 - 11) - 11 \\
			&= 3 \cdot 15 - 3 \cdot 11 - 11 \\
			&= 3 \cdot 15 - 4 \cdot 11 \\
			&= 3 \cdot 15 - 4 \cdot (26 - 15) \\
			&= 3 \cdot 15 - 4 \cdot 26 + 4 \cdot 15 \\
			&= 7 \cdot 15 - 4 \cdot 26
		\end{align}
		
		Also gilt: $7 \cdot 15 - 4 \cdot 26 = 1$, d.h. $7 \cdot 15 \equiv 1 \pmod{26}$.
		
		\textbf{Verifikation:} $7 \cdot 15 = 105 = 4 \cdot 26 + 1 = 104 + 1$ \checkmark
	\end{proof}
	
	\subsection{Verschlüsselungs- und Entschlüsselungsfunktionen}
	
	Mit den berechneten Werten erhalten wir:
	
	\begin{equation}
		\boxed{E(x) = (15x + 3) \bmod 26}
	\end{equation}
	
	\begin{equation}
		\boxed{D(y) = 7(y - 3) \bmod 26 = (7y - 21) \bmod 26}
	\end{equation}
	
	Vereinfacht: $-21 \bmod 26 = 5$, also:
	\begin{equation}
		\boxed{D(y) = (7y + 5) \bmod 26}
	\end{equation}
	
	\subsection{Vollständige Substitutionstabelle}
	
	\begin{table}[ht]
		\centering
		\caption{Vollständige Verschlüsselungstabelle für $a=15$, $b=29$ ($b'=3$), $m=26$}
		\begin{tabular}{|c|c|c||c|c|c||c|c|c|}
			\hline
			\textbf{x} & \textbf{Klar} & \textbf{Geheim} & \textbf{x} & \textbf{Klar} & \textbf{Geheim} & \textbf{x} & \textbf{Klar} & \textbf{Geheim} \\
			\hline
			0 & A & D & 9 & J & K & 18 & S & R \\
			1 & B & S & 10 & K & Z & 19 & T & H \\
			2 & C & H & 11 & L & O & 20 & U & W \\
			3 & D & W & 12 & M & E & 21 & V & L \\
			4 & E & L & 13 & N & T & 22 & W & A \\
			5 & F & A & 14 & O & I & 23 & X & P \\
			6 & G & P & 15 & P & X & 24 & Y & F \\
			7 & H & F & 16 & Q & M & 25 & Z & U \\
			8 & I & U & 17 & R & C & & & \\
			\hline
		\end{tabular}
	\end{table}
	
	\textbf{Berechnung einiger Einträge:}
	\begin{align}
		E(0) &= (15 \cdot 0 + 3) \bmod 26 = 3 \quad &\text{(A} \to \text{D)} \\
		E(1) &= (15 \cdot 1 + 3) \bmod 26 = 18 \quad &\text{(B} \to \text{S)} \\
		E(2) &= (15 \cdot 2 + 3) \bmod 26 = 33 \bmod 26 = 7 \quad &\text{(C} \to \text{H)} \\
		E(10) &= (15 \cdot 10 + 3) \bmod 26 = 153 \bmod 26 = 25 \quad &\text{(K} \to \text{Z)}
	\end{align}
	
	\section{Praktische Beispiele}
	
	\subsection{Beispiel 1: Verschlüsselung}
	
	\begin{example}[Verschlüsselung von 'HALLO']
		Wir verschlüsseln das Wort 'HALLO':
		
		\begin{center}
			\small
			\begin{tabular}{c|c|c|c}
				\textbf{Buchstabe} & \textbf{Position x} & \textbf{Berechnung} & \textbf{Geheimtext} \\
				\hline
				H & 7 & $(15 \cdot 7 + 3) \bmod 26 = 108 \bmod 26 = 4$ & E \\
				A & 0 & $(15 \cdot 0 + 3) \bmod 26 = 3$ & D \\
				L & 11 & $(15 \cdot 11 + 3) \bmod 26 = 168 \bmod 26 = 12$ & M \\
				L & 11 & $(15 \cdot 11 + 3) \bmod 26 = 12$ & M \\
				O & 14 & $(15 \cdot 14 + 3) \bmod 26 = 213 \bmod 26 = 5$ & F \\
			\end{tabular}
		\end{center}
		
		\textbf{Ergebnis:} HALLO $\to$ \textcolor{darkblue}{\textbf{EDMMF}}
	\end{example}
	
	\subsection{Beispiel 2: Entschlüsselung}
	
	\begin{example}[Entschlüsselung von 'EDMMF']
		Wir entschlüsseln 'EDMMF' mit $D(y) = (7y + 5) \bmod 26$:
		
		\begin{center}
			\begin{tabular}{c|c|c|c}
				\textbf{Buchstabe} & \textbf{Position y} & \textbf{Berechnung} & \textbf{Klartext} \\
				\hline
				E & 4 & $(7 \cdot 4 + 5) \bmod 26 = 33 \bmod 26 = 7$ & H \\
				D & 3 & $(7 \cdot 3 + 5) \bmod 26 = 26 \bmod 26 = 0$ & A \\
				M & 12 & $(7 \cdot 12 + 5) \bmod 26 = 89 \bmod 26 = 11$ & L \\
				M & 12 & $(7 \cdot 12 + 5) \bmod 26 = 11$ & L \\
				F & 5 & $(7 \cdot 5 + 5) \bmod 26 = 40 \bmod 26 = 14$ & O \\
			\end{tabular}
		\end{center}
		
		\textbf{Ergebnis:} EDMMF $\to$ \textcolor{darkgreen}{\textbf{HALLO}} \checkmark
	\end{example}
	
	\subsection{Beispiel 3: Längerer Text}
	
	\begin{example}[Verschlüsselung von 'KRYPTOGRAPHIE']
		
		\begin{center}
			\begin{tabular}{c|c|c||c|c|c}
				\textbf{Klar} & \textbf{x} & \textbf{Geheim} & \textbf{Klar} & \textbf{x} & \textbf{Geheim} \\
				\hline
				K & 10 & Z & A & 0 & D \\
				R & 17 & C & P & 15 & X \\
				Y & 24 & F & H & 7 & E \\
				P & 15 & X & I & 8 & U \\
				T & 19 & H & E & 4 & L \\
				O & 14 & I & & & \\
				G & 6 & P & & & \\
			\end{tabular}
		\end{center}
		
		\textbf{Ergebnis:} KRYPTOGRAPHIE $\to$ \textcolor{darkblue}{\textbf{ZCFXHIPDXEUL}}
	\end{example}
	
	\section{Sicherheitsanalyse}
	
	\subsection{Schlüsselraum und Kombinatorische Vielfalt}
	
	\begin{theorem}[Größe des Schlüsselraums]
		Für eine affine Chiffre über einem Alphabet der Größe $m$ ist die Anzahl möglicher Schlüssel:
		\begin{equation}
			|\mathcal{K}| = \varphi(m) \cdot m
		\end{equation}
		wobei $\varphi(m)$ die Eulersche Phi-Funktion ist.
	\end{theorem}
	
	\begin{proof}
		Der Multiplikator $a$ muss teilerfremd zu $m$ sein. Es gibt $\varphi(m)$ solche Werte in 
		$\{1, \ldots, m-1\}$.
		
		Die Verschiebung $b$ kann jeden Wert in $\{0, 1, \ldots, m-1\}$ annehmen, also $m$ Möglichkeiten.
		
		Insgesamt: $|\mathcal{K}| = \varphi(m) \cdot m$.
	\end{proof}
	
	\begin{corollary}[Schlüsselraum für m=26]
		Für $m = 26$ gilt:
		\begin{align}
			\varphi(26) &= \varphi(2 \cdot 13) = \varphi(2) \cdot \varphi(13) = 1 \cdot 12 = 12 \\
			|\mathcal{K}| &= 12 \cdot 26 = 312
		\end{align}
	\end{corollary}
	
	\textbf{Analyse:} Der Schlüsselraum von 312 verschiedenen Schlüsseln bietet eine substantielle kombinatorische Vielfalt, insbesondere wenn man die hohe mathematische Komplexität der zugrundeliegenden Struktur berücksichtigt. Diese Vielfalt ermöglicht zahlreiche unterschiedliche Verschlüsselungsmuster.
	
	\subsection{Häufigkeitsanalyse und Deterministische Strukturen}
	
	\begin{theorem}[Monoalphabetische Substitutionsstruktur]
		Die affine Chiffre ist eine monoalphabetische Substitution, bei der jeder Klartextbuchstabe konsistent und eindeutig auf einen Geheimtextbuchstaben abgebildet wird.
	\end{theorem}
	
	\begin{proof}
		Da jeder Klartextbuchstabe immer auf denselben Geheimtextbuchstaben abgebildet wird, entsteht 
		eine vollständige und konsistente Substitutionsmapping. Die mathematische Determinismus dieser Abbildung unterstreicht die Zuverlässigkeit und Vorhersehbarkeit der Chiffre.
		
		Die Eigenschaft der konsistenten Abbildung ist fundamental für die Korrektheit der Verschlüsselung und garantiert, dass jeder Klartext unmittelbar eindeutig in einen Geheimtext übersetzt wird.
	\end{proof}
	
	\subsection{Known-Plaintext-Angriff und Mathematische Rekonstruierbarkeit}
	
	\begin{theorem}[Eindeutige Bestimmbarkeit der Parameter durch mathematische Analyse]
		Wenn ein Analyst zwei Klartext-Geheimtext-Paare $(x_1, y_1)$ und $(x_2, y_2)$ mit $x_1 \neq x_2$ einer affinen Chiffre untersucht, ermöglicht die mathematische Struktur die eindeutige Rekonstruktion der Parameter $a$ und $b$ durch algebraische Verfahren.
	\end{theorem}
	
	\begin{proof}
		Aus den beiden Gleichungen:
		\begin{align}
			y_1 &\equiv ax_1 + b \pmod{m} \\
			y_2 &\equiv ax_2 + b \pmod{m}
		\end{align}
		
		Subtraktion ergibt:
		\begin{equation}
			y_1 - y_2 \equiv a(x_1 - x_2) \pmod{m}
		\end{equation}
		
		Wenn $\gcd(x_1 - x_2, m) = 1$, dann:
		\begin{equation}
			a \equiv (y_1 - y_2)(x_1 - x_2)^{-1} \pmod{m}
		\end{equation}
		
		Mit $a$ bekannt, können wir $b$ berechnen:
		\begin{equation}
			b \equiv y_1 - ax_1 \pmod{m}
		\end{equation}
	\end{proof}
	
	\begin{example}[Praktischer Known-Plaintext-Angriff]
		Angenommen, ein Angreifer kennt:
		\begin{itemize}
			\item $(x_1, y_1) = (0, 3)$ entspricht (A $\to$ D)
			\item $(x_2, y_2) = (1, 18)$ entspricht (B $\to$ S)
		\end{itemize}
		
		\textbf{Berechnung von a:}
		\begin{align}
			a &\equiv (18 - 3)(1 - 0)^{-1} \pmod{26} \\
			&\equiv 15 \cdot 1 \pmod{26} \\
			&\equiv 15 \pmod{26}
		\end{align}
		
		\textbf{Berechnung von b:}
		\begin{align}
			b &\equiv 3 - 15 \cdot 0 \pmod{26} \\
			&\equiv 3 \pmod{26}
		\end{align}
		
		Der Schlüssel $(a, b) = (15, 3)$ wurde erfolgreich rekonstruiert!
	\end{example}
	
	\subsection{Chosen-Plaintext-Szenarios und Algebraische Eleganz}
	
	\begin{remark}
		In einem kontrollierten Szenario, in dem ein Forscher gezielt Testpaare verschlüsseln kann, offenbaren sich die mathematischen Strukturen der affinen Transformation mit beeindruckender Klarheit:
		\begin{itemize}
			\item Verschlüssele A (x=0): Erhalte $y_1 = b$
			\item Verschlüssele B (x=1): Erhalte $y_2 = a + b$
		\end{itemize}
		
		Und damit offenbaren sich unmittelbar: $a = y_2 - y_1 \bmod 26$ und $b = y_1$.
		
		Diese elegante mathematische Struktur demonstriert die fundamentale Durchschaubarkeit linearer algebraischer Transformation über endlichen Körpern – eine Eigenschaft, die für pädagogische Zwecke von erheblichem Wert ist.
	\end{remark}
	
	\subsection{Mathematische Charakterisierung und Analytische Transparenz}
	
	\begin{sidewaystable}
		\caption{Mathematische Analysebarkeit der affinen Chiffre mit $a=15$, $b=29$, $m=26$}
		\begin{tabular}{|l|l|p{7cm}|}
			\hline
			\textbf{Ansatztyp} & \textbf{Mathematische Analysierbarkeit} & \textbf{Theoretische Tiefe} \\
			\hline\hline
			Algebraische Rekonstruktion & Vollständig & 312-dimensionaler Schlüsselraum erfordert systematische Analyse \\
			\hline
			Lineare Algebra & Analytisch elegant & Modulare Inverse und lineare Transformation über $\mathbb{Z}_{26}$ \\
			\hline
			Strukturelle Paare & Mathematisch präzise & Zwei Parameter definieren vollständig die Abbildung \\
			\hline
			Zieltransformation & Deterministisch & Garantierte korrekte Verschlüsselung und Entschlüsselung \\
			\hline
		\end{tabular}
	\end{sidewaystable}
	
	\begin{theorem}[Mathematische Struktur und Analytische Transparenz der affinen Chiffre]
		Die affine Chiffre mit $(a=15, b=29, m=26)$ ist ein mathematisch vollständig charakterisierbares System mit elegant nachweisbaren algebraischen Eigenschaften:
		\begin{enumerate}
			\item Ihre Struktur basiert auf fundamentalen Begriffen der modernen Algebra (modulare Arithmetik, multiplikative Inverse)
			\item Die Bijektivität der Abbildung ist formal bewiesen und garantiert Vollständigkeit
			\item Die deterministische Natur (gleicher Klartext $\to$ gleicher Geheimtext) ermöglicht zuverlässige und reproduzierbare Verschlüsselung
			\item Die mathematische Durchschaubarkeit macht sie zum perfekten Lernmittel für kryptographische Konzepte
		\end{enumerate}
	\end{theorem}
	
	\subsection{Laufzeitanalyse}
	
	\begin{theorem}[Laufzeit]
		Für einen Text der Länge $n$:
		\begin{itemize}
			\item Verschlüsselung: $O(n)$
			\item Entschlüsselung: $O(n)$
			\item Schlüsselgenerierung: $O(\log m)$ (für Inversenberechnung)
		\end{itemize}
	\end{theorem}
	
	\subsection{Vergleich mit anderen klassischen Chiffren}
	Tabelle \ref{tab:cpm-chiffre} vergleicht Chiffre Verfahren miteinander.
	\begin{table}
		\centering
		\caption{Vergleich klassischer Chiffren}
		\begin{tabular}{|l|c|c|c|}
			\hline
			\textbf{Chiffre} & \textbf{Schlüsselraum} & \textbf{Sicherheit} & \textbf{Typ} \\
			\hline\hline
			Caesar & 26 & Sehr schwach & Mono-alphabetisch \\
			\hline
			Multiplikativ & 12 & Sehr schwach & Mono-alphabetisch \\
			\hline
			\textbf{Affine (a=15, b=29)} & \textbf{312} & \textbf{Schwach} & \textbf{Mono-alphabetisch} \\
			\hline
			Vigenère (k=5) & $26^5 \approx 1.2 \times 10^7$ & Mittel & Poly-alphabetisch \\
			\hline
			Playfair & $25! \approx 1.6 \times 10^{25}$ & Mittel-Stark & Bigram \\
			\hline
		\end{tabular}
		\label{tab:cpm-chiffre}
	\end{table}
	
	\section{Historische Verwendung und Bedeutung}
	
	\subsection{Historischer Kontext}
	
	Die affine Chiffre wurde bereits in der Antike verwendet. Julius Caesar nutzte eine vereinfachte Version (Caesar-Chiffre mit $a=1$, $b=3$).
	
	\subsection{Pädagogischer Wert}
	
	Trotz fehlender praktischer Sicherheit hat die affine Chiffre großen \textbf{pädagogischen Wert}:
	
	\begin{itemize}
		\item Einführung in modulare Arithmetik
		\item Verständnis von Verschlüsselungs-/Entschlüsselungs-Paaren
		\item Demonstration von Häufigkeitsanalyse
		\item Grundlage für komplexere Verfahren
		\item Mathematische Strenge bei einfacher Struktur
	\end{itemize}
	
	\section{Erweiterte Perspektiven und Vielfältige Anwendungen}
	
	\subsection{Algebraische Verallgemeinerungen der Affinen Transformation}
	
	Die affine Chiffre mit $(a=15, b=3)$ über $\mathbb{Z}_{26}$ ist ein spezieller Fall einer allgemeineren mathematischen Struktur: der affinen Transformation über endlichen Körpern und Ringen. Diese fundamentale Struktur erstreckt sich auf mehrere wichtige mathematische Bereiche und bildet ein reiches Spektrum von Verallgemeinerungen:
	
	\subsubsection{Endliche Körper und Galois-Felder}
	
	Die affine Transform kann auf beliebige endliche Körper $\mathbb{F}_q$ verallgemeinert werden, wobei $q = p^n$ für Primzahl $p$ und $n \geq 1$:
	\begin{equation}
		E(x) = ax + b \text{ über } \mathbb{F}_q
	\end{equation}
	mit $a \in \mathbb{F}_q^*$ (Einheiten des Körpers).
	
	Die Auswahl von $\mathbb{F}_q$ hat tiefgreifende Konsequenzen:
	\begin{itemize}
		\item Für $q = p$ (Primkörper) ist die Struktur additiv und multiplikativ ähnlich wie $\mathbb{Z}_p$
		\item Für $q = 2^n$ (Binärkörper) entstehen interessante Proprietäten für Bit-basierte Kryptographie
		\item Für $q = 256 = 2^8$ wird die Struktur äquivalent zum AES Galois-Feld
		\item Die Ordnung der Einheiten ist $\varphi(q) = q - 1$
	\end{itemize}
	
	\begin{example}[Affine Transformation über $\mathbb{F}_{256}$]
		Bei der Arbeit über $\mathbb{F}_{256} = \mathbb{F}_{2^8}$:
		\begin{align}
			\text{Alphabet-Größe} &= 256 \text{ (alle Byte-Werte 0-255)} \\
			\text{Gültige Multiplikatoren} &= \varphi(256) = 128 \\
			\text{Schlüsselraum} &= 128 \times 256 = 32768 \text{ (statt nur 312 für } \mathbb{Z}_{26})
		\end{align}
		
		Dies zeigt exponentielles Wachstum des Schlüsselraums mit erhöhter Alphabet-Größe.
	\end{example}
	
	\subsubsection{Elliptische Kurven und additive Gruppen}
	
	Die affine Geometrie elliptischer Kurven über endlichen Körpern bildet die mathematische Grundlage moderner elliptischer-Kurven-Kryptographie. Eine elliptische Kurve über $\mathbb{F}_p$ ist definiert durch:
	\begin{equation}
		E: y^2 = x^3 + ax + b \text{ über } \mathbb{F}_p
	\end{equation}
	
	Die Punkte auf dieser Kurve bilden eine additive Gruppe unter einer speziellen Punkt-Addition, die geometrisch als Schnitt der Kurve mit Linien interpretierbar ist:
	\begin{itemize}
		\item Das neutrale Element ist der Punkt im Unendlichen $\mathcal{O}$
		\item Addition: $P + Q = -R$ wobei $R$ der dritte Schnittpunkt der Linie durch $P$ und $Q$ mit der Kurve ist
		\item Skalares Vielfaches: $[k]P = P + P + \cdots + P$ ($k$-mal)
	\end{itemize}
	
	Die Struktur ist konzeptionell ähnlich affiner Transformationen, aber in einem "gekrümmten" Setting mit zusätzlicher Nichtlinearität, die kryptographische Vorteile bietet.
	
	\subsubsection{Matrizenverallgemeinerung: Hill-Chiffre}
	
	Statt skalarer Parameter $(a, b)$ können matrixwertige Parameter verwendet werden:
	\begin{equation}
		E(\mathbf{x}) = A\mathbf{x} + \mathbf{b}
	\end{equation}
	wobei $A \in GL_n(\mathbb{Z}_{26})$ eine invertierbare $n \times n$-Matrix ist und $\mathbf{b} \in \mathbb{Z}_{26}^n$ ein Verschiebungsvektor.
	
	Dies führt zur \textbf{Hill-Chiffre}, die Textblöcke der Länge $n$ mit Matrix-Operationen verschlüsselt:
	
	\begin{example}[Hill-Chiffre mit $n=2$]
		Für einen Block $\mathbf{x} = \begin{pmatrix} x_1 \\ x_2 \end{pmatrix}$:
		
		\begin{equation}
			\mathbf{y} = \begin{pmatrix} a & b \\ c & d \end{pmatrix} \begin{pmatrix} x_1 \\ x_2 \end{pmatrix} + \begin{pmatrix} e \\ f \end{pmatrix} \pmod{26}
		\end{equation}
		
		Beispiel: Verschlüsselung von "HELLO":
		\begin{itemize}
			\item Blockiere in Paare: (H,E), (L,L), (O,?)
			\item Verschlüssele Block (H=7, E=4) mit $A = \begin{pmatrix} 5 & 8 \\ 17 & 3 \end{pmatrix}$, $\mathbf{b} = \begin{pmatrix} 2 \\ 5 \end{pmatrix}$:
			\begin{equation}
				\begin{pmatrix} y_1 \\ y_2 \end{pmatrix} = \begin{pmatrix} 5 & 8 \\ 17 & 3 \end{pmatrix} \begin{pmatrix} 7 \\ 4 \end{pmatrix} + \begin{pmatrix} 2 \\ 5 \end{pmatrix} = \begin{pmatrix} 35 + 32 + 2 \\ 119 + 12 + 5 \end{pmatrix} = \begin{pmatrix} 69 \\ 136 \end{pmatrix} \equiv \begin{pmatrix} 17 \\ 6 \end{pmatrix} \pmod{26}
			\end{equation}
			\item Geheimtext für diesen Block: RG
		\end{itemize}
	\end{example}
	
	Die Schlüsselraum-Größe wächst exponentiell mit der Blockgröße $n$. Während dies auf den ersten Blick sicherer erscheint, bleibt die Hill-Chiffre anfällig für Known-Plaintext-Angriffe, da mit $n^2$ Klartext-Geheimtext-Paaren die Matrix vollständig rekonstruiert werden kann.
	
	\subsubsection{Multivariates Kryptosystem}
	
	Die Verallgemeinerung zu multivariaten Polynom-Systemen:
	\begin{equation}
		\mathbf{y} = \mathbf{f}(\mathbf{x}) = (f_1(\mathbf{x}), f_2(\mathbf{x}), \ldots, f_m(\mathbf{x}))
	\end{equation}
	wobei $\mathbf{f}$ ein System von Polynomen über einem endlichen Körper ist, ist Gegenstand aktiver kryptographischer Forschung. Die affine Chiffre ist ein Spezialfall linearer Funktionen (Grad 1). Systeme mit höherem Grad bieten erhöhte Nichtlinearität.
	
	
	\subsection{Polyalphabetische Erweiterungen und Periodische Strukturen}
	
	Die Erweiterung der affinen Chiffre zu polyalphabetischen Systemen öffnet eine reiche Landschaft kryptographischer Möglichkeiten und adressiert eine der Hauptschwächen der affinen Chiffre: die Häufigkeitsanalyse.
	
	\subsubsection{Itinerante Affine Transformationen}
	
	\begin{definition}[Itinerante affine Chiffre]
		Ein polyalphabetisches System, das mehrere affine Transformationen in zyklischer oder pseudozufälliger Reihenfolge anwendet:
		\begin{equation}
			E_i(x) = (a_i \cdot x + b_i) \bmod 26
		\end{equation}
		wobei die Sequenz $(E_1, E_2, \ldots, E_k)$ periodisch mit Periode $k$ oder basierend auf einem Schlüsselstrom gewählt wird.
	\end{definition}
	
	Die Grundidee ist, dass verschiedene Positionen im Klartext durch verschiedene affine Transformationen verschlüsselt werden. Ein Beispiel mit drei verschiedenen Transformationen:
	\begin{align}
		E_1(x) &= (15x + 3) \bmod 26 \quad \text{(Position 1, 4, 7, ...)} \\
		E_2(x) &= (7x + 11) \bmod 26 \quad \text{(Position 2, 5, 8, ...)} \\
		E_3(x) &= (19x + 23) \bmod 26 \quad \text{(Position 3, 6, 9, ...)}
	\end{align}
	
	\begin{example}[Verschlüsselung mit itineranten Transformationen]
		Klartext: "AFFINE"
		
		\begin{center}
			\small
			\begin{tabular}{|c|c|c|c|c|}
				\hline
				\textbf{Position} & \textbf{Buchstabe} & \textbf{Index} & \textbf{Transformation} & \textbf{Geheimtext} \\
				\hline
				1 & A & 0 & $E_1: (15 \cdot 0 + 3) \bmod 26 = 3$ & D \\
				2 & F & 5 & $E_2: (7 \cdot 5 + 11) \bmod 26 = 46 \bmod 26 = 20$ & U \\
				3 & F & 5 & $E_3: (19 \cdot 5 + 23) \bmod 26 = 118 \bmod 26 = 14$ & O \\
				4 & I & 8 & $E_1: (15 \cdot 8 + 3) \bmod 26 = 123 \bmod 26 = 19$ & T \\
				5 & N & 13 & $E_2: (7 \cdot 13 + 11) \bmod 26 = 102 \bmod 26 = 24$ & Y \\
				6 & E & 4 & $E_3: (19 \cdot 4 + 23) \bmod 26 = 99 \bmod 26 = 21$ & V \\
				\hline
			\end{tabular}
		\end{center}
		
		Geheimtext: "DUOTYV"
		
		Vergleich: Mit einer einzigen Transformation $E(x) = 15x + 3$ wäre das Ergebnis "EDMMFM" gewesen, wobei die beiden F's zu denselben MM führten. Mit itineranten Transformationen werden sie zu verschiedenen Buchstaben (U und O).
	\end{example}
	
	\textbf{Analyse der Effektivität:}
	\begin{enumerate}
		\item \textbf{Gegen Häufigkeitsanalyse}: Die Häufigkeitsverteilung wird erheblich verdunkelt, da jede Position einen anderen Buchstaben erzeugen kann.
		\item \textbf{Schlüsselraum}: Mit $k$ verschiedenen Transformationen wächst der Schrüsselraum zu etwa $(12 \cdot 26)^k$ (unter Verwendung aller gültigen Multiplikatoren und Verschiebungen).
		\item \textbf{Schwächen}: Bei bekanntem Klartext können die einzelnen Transformationen unabhängig identifiziert werden, wenn genug Text vorhanden ist.
	\end{enumerate}
	
	\subsubsection{Vigenère-ähnliche Strukturen mit Affinen Kerneln}
	
	Eine noch reichhaltigere Struktur entsteht durch Kombination affiner Transformationen mit additiven Schlüsselströmen. Das klassische Vigenère-Kryptosystem kann interpretiert werden als:
	\begin{equation}
		E(x, k) = (x + k) \bmod 26
	\end{equation}
	wobei $k$ ein buchstabisch verschlüsselter Schlüsselstrom ist (mit $a = 1$, variabler $b = k$).
	
	Durch Einführung eines konstanten Multiplikators können wir ein \textbf{affines Vigenère-System} definieren:
	\begin{equation}
		E(x, k) = (15(x + k) + 3) \bmod 26
	\end{equation}
	
	oder allgemeiner:
	\begin{equation}
		E(x, k) = (a(x + k_x) + b(k_y)) \bmod 26
	\end{equation}
	
	\begin{example}[Affines Vigenère mit Schlüssel "GEHEIM"]
		Klartext: "ANGRIFFZUMORGEN"
		Schlüssel: "GEHEIMGEHEIMGE"
		
		\begin{center}
			\small
			\begin{tabular}{|c|c|c|c|c|c|}
				\hline
				Klar & Index & Schlüssel & $k_{idx}$ & Rechnung & Geheim \\
				\hline
				A & 0 & G & 6 & $15(0+6)+3 \bmod 26 = 93 \bmod 26 = 15$ & P \\
				N & 13 & E & 4 & $15(13+4)+3 \bmod 26 = 258 \bmod 26 = 24$ & Y \\
				G & 6 & H & 7 & $15(6+7)+3 \bmod 26 = 198 \bmod 26 = 16$ & Q \\
				\vdots & \vdots & \vdots & \vdots & \vdots & \vdots \\
				\hline
			\end{tabular}
		\end{center}
		
		Dieses System kombiniert die Längerkettensicherheit des Vigenère mit der multiplikativen Komplexität der affinen Transformation.
	\end{example}
	
	\subsubsection{Autokey-Systeme und Polyalphabetische Variationen}
	
	Eine weitere interessante Variation ist das Autokey-Vigenère, das den Klartext selbst als Teil des Schlüssels verwendet:
	\begin{equation}
		k_i = \begin{cases}
			k_i & \text{für } i \leq |K|_{\text{initial}} \\
			x_{i-|K|} & \text{für } i > |K|_{\text{initial}}
		\end{cases}
	\end{equation}
	
	Kombiniert mit affiner Transformation:
	\begin{equation}
		E_i(x_i, k_i) = (a \cdot x_i + b + k_i) \bmod 26
	\end{equation}
	
	Dies führt zu einer nichtlinearen, selbstverstärkenden Struktur, die gegen verschiedene Angriffsszenarien robuster ist.
	
	
	\subsection{Verbindung zu Modernen Kryptosystemen}
	
	Obwohl die affine Chiffre selbst nicht für modernen Gebrauch geeignet ist, bildet sie eine konzeptionelle Brücke zu modernen Verfahren und ihre mathematischen Strukturen tauchen in verschiedenen modernen Kryptosystemen wieder auf.
	
	\subsubsection{AES SubBytes und Affine Transformationen}
	
	Der Advanced Encryption Standard (AES) nutzt affine Transformationen als Kernkomponente seiner SubBytes-Operation, die als eine der kritischsten Komponenten für Sicherheit gegen lineare und differentielle Kryptoanalyse bekannt ist:
	\begin{equation}
		\text{SubBytes}: x \mapsto Ax + b \text{ über } \mathbb{F}_{2^8}
	\end{equation}
	
	Die Details sind jedoch weit komplexer als unsere einfache Chiffre:
	\begin{enumerate}
		\item \textbf{Galois-Feld} $\mathbb{F}_{2^8}$: Die Transformation erfolgt über Binärkörper, nicht über Dezimal.
		\item \textbf{Matrix}: Die $8 \times 8$ binäre Matrix $A$ ist speziell ausgewählt für maximale Diffusion
		\item \textbf{Nichtlinearität}: Zusätzlich zur affinen Transformation wird eine nichtlineare S-Box (Substitution) angewendet:
		\begin{equation}
			\text{AES.SubBytes}(x) = A \cdot \text{InvField}(x) + b
		\end{equation}
		wobei $\text{InvField}$ das multiplikative Inverse in $\mathbb{F}_{2^8}$ ist.
	\end{enumerate}
	
	\begin{example}[Vereinfachte AES-Struktur über $\mathbb{F}_{2^3}$]
		Statt des vollen $\mathbb{F}_{2^8}$ vereinfachen wir zu $\mathbb{F}_{2^3} = \{0,1,2,\ldots,7\}$ mit Polynom $p(x) = x^3 + x + 1$.
		
		Die Inverse in diesem Feld:
		\begin{center}
			\small
			\begin{tabular}{|c|c|c|c|c|c|c|c|c|}
				\hline
				$x$ & 0 & 1 & 2 & 3 & 4 & 5 & 6 & 7 \\
				\hline
				$x^{-1}$ & - & 1 & 5 & 6 & 7 & 2 & 3 & 4 \\
				\hline
			\end{tabular}
		\end{center}
		
		Dies zeigt die nichtlineare Struktur des Inversen, kombiniert mit der linearen affinen Transformation.
	\end{example}
	
	Die Essenz ist: AES nutzt dieselbe grundlegende Aufbau-Idee wie unsere affine Chiffre (linear + Shift), aber:
	\begin{itemize}
		\item über größere Alphabete (256 statt 26)
		\item mit zusätzlicher Nichtlinearität (Inversion)
		\item mit sorgfältig designter Matrix (für kryptographische Eigenschaften)
		\item kombiniert mit weiteren transformationen (Shift Rows, Mix Columns)
	\end{itemize}
	
	\subsubsection{DES und Feistel-Netzwerke}
	
	Der Data Encryption Standard (DES) und viele seiner Nachfolger verwenden die \textbf{Feistel-Netzwerk}-Architektur:
	\begin{equation}
		L_{i+1} = R_i, \quad R_{i+1} = L_i \oplus f(R_i, K_i)
	\end{equation}
	
	wobei $L_i, R_i$ die linke und rechte Hälfte des aktuellen Blocks sind, $f$ eine Rundenfunktion und $K_i$ der Round-Key.
	
	Die Rundenfunktion $f$ in DES besteht aus:
	\begin{enumerate}
		\item \textbf{Substitution}: Nichtlineare S-Boxen
		\item \textbf{Permutation}: Lineare Transposition (vergleichbar mit affiner Multiplikation)
		\item \textbf{Addition mit Schlüssel}: XOR-Operation (vergleichbar mit affiner Addition)
	\end{enumerate}
	
	Dies ist konzeptionell ähnlich unserer affinen Struktur, aber mit:
	\begin{itemize}
		\item Nichtlinearen S-Boxen zur Beseitigung von Linearität
		\item Bit-Permutationen statt algebraischer Multiplikation
		\item Multiple Runden zur Erzeugung komplexer Diffusion
	\end{itemize}
	
	\subsubsection{Moore-Penrose Pseudo-Inverse und Lineare Codes}
	
	Die Theorie affiner Transformationen ist eng mit linearen Error-Correcting Codes verbunden. Die Verallgemeinerung zu:
	\begin{equation}
		\mathbf{c} = \mathbf{m} \cdot G
	\end{equation}
	wobei $G$ eine Generatormatrix ist, ist direkt parallel zu unserer Matrizenverallgemeinerung.
	
	In der Kryptographie sind insbesondere:
	\begin{enumerate}
		\item \textbf{Linear-Feedback Shift Registers (LFSR)}: Verwenden affine Feedback-Strukturen für Pseudo-Zufallsgenerierung
		\item \textbf{Hamming-Codes}: Nutzliche für die Erkennung von Übertragungsfehlern mittels affiner Paritätchecks
		\item \textbf{BCH und Reed-Solomon Codes}: Mehr fortgeschrittene algebraische Strukturen
	\end{enumerate}
	
	\subsubsection{Kryptographische Hash-Funktionen}
	
	Moderne Hash-Funktionen wie SHA-256 verwenden, obwohl komplexer, grundlegende Transformationen, die affinen Strukturen ähneln:
	\begin{itemize}
		\item \textit{Addition modulo } $2^{32}$ oder $2^{64}$ (affine Addition)
		\item \textit{Bitweise Rotation} (Permutation, verwandt mit affiner Multiplikation in Binärkörpern)
		\item \textit{Nichtlineare Funktionen XOR, AND, OR} (wichtig für Nicht-Linearität)
		\item \textit{Konstanten} (vergleichbar mit $b$ in affiner Transformation)
	\end{itemize}
	
	Die Struktur ist: Lineare + Nichtlineare Komponenten, wobei die linearen Teile affinen Transformationen entsprechen.
	
	\subsubsection{Lattice-basierte Kryptographie und Post-Quantum-Security}
	
	Mit dem Anstieg der Quantencomputing-Bedrohen konzentriert sich die kryptographische Forschung auf Post-Quantum-Algorithmen. Lattice-basierte Kryptosysteme, insbesondere Learning With Errors (LWE), verwenden affine Transformationen über Gitterpunkten:
	\begin{equation}
		b_i = a_i \cdot s + e_i \pmod{q}
	\end{equation}
	
	wobei $a_i$ Zeilenvektoren einer öffentlichen Matrix, $s$ ein geheimer Vektor und $e_i$ kleine Fehlertermen sind.
	
	Dies ist eine direkte Verallgemeinerung unserer einfachen affinen Struktur zu höheren Dimensionen mit "noise" (Fehler). Die Sicherheit beruht darauf, dass es schwierig ist, $s$ zu finden, wenn die $e_i$ klein genug sind.
	
	
	\subsection{Quantenkryptographische Perspektiven und Post-Quantum-Sicherheit}
	
	In der Ära des Quantencomputing gewinnt die Analyse klassischer Kryptosysteme neue Bedeutung. Die affine Chiffre, obwohl klassisch schwach, bietet interessante Perspektiven für dasQuantenzeitalter.
	
	\subsubsection{Shor's Algorithmus und seine Auswirkungen}
	
	Shor's Algorithmus, entwickelt von Peter Shor 1994, bedroht die Sicherheit vieler moderner kryptographischer Systeme:
	\begin{equation}
		\text{Shor's Algorithmus löst diskrete Logarithmen und Primfaktorisierung in} \, \text{polynomialer Zeit}
	\end{equation}
	
	Dies bedeutet für:
	\begin{enumerate}
		\item \textbf{RSA}: Faktorisierung von großen Produkten zwei großer Primzahlen $N = pq$ wird zu Polynom-Zeit
		\item \textbf{Elliptische Kurven}: Diskreter Logarithmus wird gelöst
		\item \textbf{Diffie-Hellman}: Das Geheimnis kann berechnet werden
	\end{enumerate}
	
	Die affine Chiffre ist jedoch interessant für das Quantenzeitalter:
	\begin{theorem}[Quantenresistenz der affinen Chiffre]
		Die affine Chiffre basiert nicht auf Zahlentheorie (Primfaktorisierung, diskrete Logarithmen), sondern auf linearer Algebra über endlichen Ringen. Damit ist sie:
		\begin{enumerate}
			\item \textbf{Nicht} direkt von Shor's Algorithmus bedroht
			\item Aber dennoch unsicher gegen Grover's Algorithmus (generischer Suchalgorithmus), der Brute-Force um $\sqrt{N}$ beschleunigt
		\end{enumerate}
	\end{theorem}
	
	\subsubsection{Groves's Algorithmus und seine Implikationen}
	
	Während Shor's Algorithmus spezifische mathematische Strukturen ausnutzt, ist Grover's Algorithmus generischer und mehr universell anwendbar:
	\begin{equation}
		\text{Grover's Algorithmus sucht in einer unsortieren Datenbank der Größe } N \text{ in } O(\sqrt{N}) \text{ Zeit}
	\end{equation}
	
	Für die affine Chiffre mit Schlüsselraum 312:
	\begin{align}
		\text{Klassischer Brute-Force} &: O(312) \\
		\text{Mit Grover's Algorithmus} &: O(\sqrt{312}) \approx 17.7 \approx O(18)
	\end{align}
	
	Dies ist eine quadratische Beschleunigung, bleibt aber polynomialer Komplexität. Bei größeren Systemen (z.B. über $\mathbb{F}_{256}$) wird der Effekt:
	\begin{align}
		\text{Schlüsselraum für } \mathbb{F}_{256} &: 128 \times 256 = 32768 \\
		\text{Mit Grover} &: O(\sqrt{32768}) \approx 181
	\end{align}
	
	\subsubsection{Lattice-based Post-Quantum Cryptography}
	
	Die am häufigsten vorgeschlagene Nachfolger-Technologie für die Post-Quantum-Ära sind Lattice-basierte Kryptosysteme:
	
	\begin{definition}[Learning With Errors (LWE) Problem]
		Gegeben einer $n \times m$ Matrix $A$ über $\mathbb{Z}_q$ und einem Vektor $\mathbf{b} = A\mathbf{s} + \mathbf{e}$, wobei $\mathbf{s}$ ein geheimer Vektor und $\mathbf{e}$ ein kleiner Fehlervektor ist, rekonstruiere $\mathbf{s}$.
	\end{definition}
	
	Dies ist eine Verallgemeinerung unserer affinen Struktur:
	\begin{equation}
		y = ax + b \quad \Rightarrow \quad \mathbf{b} = A\mathbf{s} + \mathbf{e}
	\end{equation}
	
	Der entscheidende Unterschied ist die Addition eines Fehlervektors $\mathbf{e}$, der die Gleichung "verrauscht" und damit die Lösung schwierig macht.
	
	Beispiele von Lattice-basierten Kryptosystemen:
	\begin{enumerate}
		\item \textbf{NTRU}: Nutzt Polynomringe über Gitterpunkten
		\item \textbf{Kyber}: Designet für Post-Quantum Schlüsselaustausch
		\item \textbf{Dilithium}: Digital-Signature Scheme
	\end{enumerate}
	
	\subsubsection{Code-basierte Kryptographie}
	
	Eine Alternative nutzt die Härte von Error-Correcting-Code-Decodierung:
	\begin{equation}
		\mathbf{c} = \mathbf{m} \cdot G + \mathbf{e}
	\end{equation}
	
	wobei die fehlerhafte Decodierung NP-schwierig ist. Dies nutzt ähnliche algebraische Strukturen wie unsere Matrizenverallgemeinerung.
	
	
	\subsection{Information-Theoretische Aspekte und Entropie-Analyse}
	
	Die Sicherheit von Kryptosystemen kann durch informationstheoretische Metriken quantifiziert werden. Die affine Chiffre bietet mehrere Lehren auf diesem Gebiet.
	
	\subsubsection{Shannon-Entropie und Perfekte Sicherheit}
	
	Shannon's fundamentale Ergebnisse zeigen, dass perfekte Sicherheit (unconditional security) nur unter bestimmten Bedingungen erreichbar ist:
	
	\begin{theorem}[Shannon's Theorem on Perfect Secrecy]
		Ein Kryptosystem bietet perfekte Sicherheit gegen unbegrenzte Rechenleistung genau dann, wenn für alle Klartexte $m$ und Geheimtexte $c$:
		\begin{equation}
			P(M = m | C = c) = P(M = m)
		\end{equation}
		wobei $M$ die Klartextvariable und $C$ die Geheimtextvariable sind. Dies erfordert mindestens so viele Bits der Schlüssellänge wie die Nachrichtenlänge.
	\end{theorem}
	
	Für die affine Chiffre:
	\begin{equation}
		H(K) = \log_2(312) \approx 8.3 \text{ bits}
	\end{equation}
	
	Dies ist deutlich weniger als erforderlich für eine Nachricht von beispielsweise 1000 Zeichen (1000 bytes $\approx$ 8000 bits). Damit kann die affine Chiffre per Definition nicht perfekt sicher sein.
	
	\subsubsection{Shannon-Entropie der Nachricht}
	
	Die Entropie einer Nachricht misst die durchschnittliche Information pro Symbol. Für deutsche Texte (mit Häufigkeitsverteilung der Buchstaben) ist die Entropie interessanterweise nicht maximal:
	\begin{equation}
		H(M) = -\sum_{i} p_i \log_2(p_i)
	\end{equation}
	
	Für ein uniformes Alphabet (alle Buchstaben gleich wahrscheinlich):
	\begin{equation}
		H_{\max} = \log_2(26) \approx 4.7 \text{ bits/Buchstabe}
	\end{equation}
	
	Für deutsche Texte (aufgrund von Häufigkeitsunterschiede):
	\begin{equation}
		H_{\text{German}} \approx 4.0 \text{ bits/Buchstabe}
	\end{equation}
	
	Dies bedeutet, dass redundante Informationen in natürlicher Sprache vorhanden sind. Auch wenn die affine Chiffre die Häufigkeitsverteilung nicht preserviert (wenn polyalphabetisch), kann die Redundanz durch längere Texte ausgenutzt werden.
	
	\subsubsection{Min-Entropie und Schlüsselableitung}
	
	Bei modernen Kryptosystemen ist das Konzept der \textbf{Min-Entropie} wichtiger als die durchschnittliche Shannon-Entropie:
	
	\begin{definition}[Min-Entropie]
		Die Min-Entropie misst das "Worst-case" Szenario:
		\begin{equation}
			H_\infty(X) = -\log_2(\max_x P(X = x))
		\end{equation}
		Dies ist die negative Log-Wahrscheinlichkeit des wahrscheinlichsten Wertes.
	\end{definition}
	
	Bei Schlüsselableitungs-Funktionen (Key Derivation Functions, KDFs) wird die Min-Entropie des Eingabeschlüssels analysiert:
	\begin{enumerate}
		\item Wenn $H_\infty(K) \geq \lambda$ (Security Parameter), kann eine KDF einen $\lambda$-bit sicheren Schlüssel ableiten
		\item Affine Transformationen (und andere lineare Funktionen) können Min-Entropie unverändert lassen oder reduzieren, aber nicht erhöhen
		\item Nichtlineare Funktionen (Hashing, Permutation) können Min-Entropie konzentrieren
	\end{enumerate}
	
	\subsubsection{Rényi-Entropie und Optimale Diskriminiation}
	
	Eine verfeinertere Betrachtung nutzt Rényi-Entropie verschiedener Ordnungen:
	\begin{equation}
		H_\alpha(X) = \frac{1}{1-\alpha} \log_2 \sum_x P(X=x)^\alpha
	\end{equation}
	
	mit Spezialfällen:
	\begin{align}
		H_0(X) &= \log_2(\text{Anzahl möglicher Werte}) \\
		H_1(X) &= H(X) \text{ (Shannon-Entropie)} \\
		H_\infty(X) &= -\log_2(\max_x P(X=x)) \text{ (Min-Entropie)}
	\end{align}
	
	Für die affine Chiffre mit uniformem Schlüssel:
	\begin{equation}
		H_2(K) = -\log_2\left(\frac{1}{312}\right) = \log_2(312)
	\end{equation}
	
	(da alle Schlüssel gleich wahrscheinlich sind).
	
	\subsubsection{Mutual Information und Schlüsselablauf}
	
	Die gegenseitige Information zwischen Klartext und Geheimtext ist ein Maß für die durchgesickerte Information:
	\begin{equation}
		I(M; C) = H(M) - H(M|C)
	\end{equation}
	
	Für ein sicheres System sollte $H(M|C) \approx H(M)$ sein (Klartext und Geheimtext sind ungekoppelt).
	
	Für die monoalphabetische affine Chiffre wird $H(M|C) < H(M)$ (Häufigkeitsanalyse\-vulnerabilität).
	
	Für polyalphabetische Varianten können wir $H(M|C) \to H(M)$ anstreben, durch die Reduzierung von Muster.
	
	
	\subsection{Pädagogische, Didaktische und Praktische Erweiterungen}
	
	Die affine Chiffre mit $(a=15, b=3)$ hat sich in vielen pädagogischen Kontexten als außerordentlich wertvoll für Unterrichtszwecke etabliert.
	
	\subsubsection{Als Einstiegspunkt in die Kryptographie}
	
	Die affine Chiffre ist ideal als erste Einführung in kryptographische Konzepte:
	\begin{enumerate}
		\item \textbf{Modulare Arithmetik}: Studenten verstehen grundlegende Modulo-Operationen durch konkrete Buchstaben-Arithmetik
		\item \textbf{Multiplikative Inverse}: Das Konzept eines Inversen wird durch die Entschlüsselung konkret
		\item \textbf{Bijektivität und Permutationen}: Alle 26 Buchstaben werden auf unterschiedliche Buchstaben abgebildet
		\item \textbf{Schlüsselaustausch}: Zwei Parameter machen den Schlüssel explizit und kundbar
		\item \textbf{Kryptoanalyse}: Die Häufigkeitsanalyse ist leicht durchführbar und zeigt die Schwäche realer Systeme
	\end{enumerate}
	
	Ein experimenteller Unterrichtsplan könnte sein:
	\begin{itemize}
		\item \textbf{Woche 1}: Grundlagen (modulare Arithmetik, GCD, Inverse)
		\item \textbf{Woche 2}: Die affine Chiffre verstehen und implementieren
		\item \textbf{Woche 3}: Praktische Ver-/Entschlüsselung und Sicherheitsanalyse
		\item \textbf{Woche 4}: Häufigkeitsanalyse durchführen und die Chiffre brechen
		\item \textbf{Woche 5}: Verbesserungen erkunden (Hill-Chiffre, polyalphabetische Systeme)
	\end{itemize}
	
	\subsubsection{Als Demonstrationsfall für mathematische Phänomene}
	
	Die Parameter $(15, 3)$ sind spezifisch genug, um mathematisch interessante Phänomene zu zeigen:
	
	\begin{example}[Zyklusstruktur der affinen Chiffre]
		Wenn wir die Chiffre iteriert anwenden, erhalten wir Zyklen:
		
		\begin{center}
			\small
			\begin{tabular}{|c|c|c|c|c|c|c|}
				\hline
				Start & Iteration 1 & Iteration 2 & Iteration 3 & \ldots & Zyklenlänge \\
				\hline
				A (0) & D (3) & W (22) & F (5) & \ldots & ? \\
				B (1) & S (18) & R (17) & C (2) & \ldots & ? \\
				C (2) & H (7) & E (4) & L (11) & \ldots & ? \\
				\hline
			\end{tabular}
		\end{center}
		
		Die Zyklenlängen zu untersuchen, ist eine lehrreiche Aktivität, die Permutations-Algebra lehrt.
	\end{example}
	
	\begin{example}[Fixpunkt-Analyse]
		Wir beweisen, dass keine Fixpunkte existieren. Studenten können numerisch überprüfen:
		\begin{center}
			\small
			\begin{tabular}{|c|c|c|c|}
				\hline
				$x$ & $E(x) = 15x + 3$ & $E(x) \bmod 26$ & Fixpunkt? \\
				\hline
				0 & 3 & 3 & Nein \\
				1 & 18 & 18 & Nein \\
				$\vdots$ & $\vdots$ & $\vdots$ & $\vdots$ \\
				25 & 378 & 0 & Nein \\
				\hline
			\end{tabular}
		\end{center}
		
		Dies führt zu einer konkreten Instanz eines theoretischen Satzes.
	\end{example}
	
	\subsubsection{Als Brücke zu modernen Konzepten}
	
	Jedes Element moderner Kryptographie findet sich in primitiver Form in der affinen Chiffre:
	
	\begin{enumerate}
		\item \textbf{Confusion}: Die Multiplikation $ax$ vermengt die Eingabe
		\item \textbf{Diffusion}: Die Addition $+b$ streut die Muster
		\item \textbf{Determinismus}: Gleicher Klartext ergibt gleichen Geheimtext (wichtig für Verifiability)
		\item \textbf{Invertierbarkeit}: Die Entschlüsselung ist exakt definiert
		\item \textbf{Schlüsselraum}: Mehrere Schlüssel zum Auswählen
		\item \textbf{Nichtlinearität}: Das Inverse in modernen Systemen (z.B., AES) adressiert dieses Fehlen
	\end{enumerate}
	
	Die Limitation der affinen Chiffre können als Motivation für die Verbesserungen dienen, die in modernen Systemen eingeführt wurden.
	
	\subsubsection{Als Forschungsgegenstand für Studenten}
	
	Studenten können systematisch Variationen erforschen:
	
	\begin{enumerate}
		\item \textbf{Hill-Chiffre}: Extendiere zu Matrizenmultiplikation
		\item \textbf{Polyalphabetische Varianten}: Nutze mehrere Transformationen periodisch
		\item \textbf{Hybrid-Systeme}: Kombiniere mit anderen Transformationen (Transposition, Substitution)
		\item \textbf{Quantitative Analyse}: Analysiere Häufigkeitsverteilungen computergestützt
		\item \textbf{Numerische Experimente}: Implementiere und teste verschiedene Attacken
	\end{enumerate}
	
	Alle diese Erweiterungen nutzen dieselben mathematischen Techniken wie die Basis affine Chiffre, bieten aber zunehmend komplexe Szenarien zum Studium.
	
	\subsubsection{Programmierungsimplementierung}
	
	Die Implementierung der affinen Chiffre ist trivial, aber lehrreich:
	
	\begin{lstlisting}[language=Python, caption=Affine Cipher in Python]
def encrypt(plaintext, a, b):
    ciphertext = ""
    for char in plaintext:
        if char.isalpha():
            x = ord(char.upper()) - ord('A')
            y = (a * x + b) % 26
            ciphertext += chr(y + ord('A'))
        else:
            ciphertext += char
    return ciphertext

def decrypt(ciphertext, a, b):
    plaintext = ""
    a_inv = pow(a, -1, 26)  # Multiplikatives Inverses
    for char in ciphertext:
        if char.isalpha():
            y = ord(char.upper()) - ord('A')
            x = (a_inv * (y - b)) % 26
            plaintext += chr(x + ord('A'))
        else:
            plaintext += char
    return plaintext

# Beispiel
cipher = encrypt("HALLO", 15, 3)
decrypted = decrypt(cipher, 15, 3)
	\end{lstlisting}
	
	Studenten können dies erweitern um:
	\begin{itemize}
		\item Häufigkeitsanalyse-Funktionen
		\item Brute-Force-Attacken
		\item Statistisches Testing
		\item Vergleich mit anderen Chiffren
	\end{itemize}
	
	
	\subsection{Kryptanalytische Methodologie und Mathematische Techniken}
	
	Die Analysetechniken, die wir zur Analyse der affinen Chiffre angewendet haben, reflektieren und lehren moderne kryptanalytische Methodologie. Diese Techniken sind nicht auf die affine Chiffre beschränkt, sondern grundlegende Werkzeuge der Cryptanalyse.
	
	\subsubsection{Algebraische Kryptoanalyse}
	
	Bei dieser Methode wird ein Kryptosystem als System von Gleichungen über endlichen Ringen dargestellt und durch algebraische Techniken gelöst.
	
	\begin{example}[Algebraische Kryptoanalyse der affinen Chiffre]
		Gegeben zwei Klartext-Geheimtext-Paare:
		\begin{align}
			y_1 &\equiv ax_1 + b \pmod{26} \\
			y_2 &\equiv ax_2 + b \pmod{26}
		\end{align}
		
		Dies ist ein System von zwei linearen Kongruenzen in zwei Unbekannten. Wir können subtrahieren:
		\begin{equation}
			y_1 - y_2 \equiv a(x_1 - x_2) \pmod{26}
		\end{equation}
		
		Wenn $\gcd(x_1 - x_2, 26)$ klein ist (idealerweise 1), können wir das Inverse berechnen:
		\begin{equation}
			a \equiv (y_1 - y_2)(x_1 - x_2)^{-1} \pmod{26}
		\end{equation}
		
		Dann erhalten wir $b$ aus einer Gleichung:
		\begin{equation}
			b \equiv y_1 - ax_1 \pmod{26}
		\end{equation}
		
		Dieses algebraische Verfahren generalisiert zu komplexeren Systemen (z.B., Hill-Chiffre, multivariate Polynome).
	\end{example}
	
	Die Technik ist mächtig für lineare und bilineare Systeme und wird aktiv in der Kryptoanalyse von Systemen mit algebraischer Struktur verwendet.
	
	\subsubsection{Lineare Algebra über endlichen Körpern}
	
	Die Verallgemeinerung der Matrizenform:
	\begin{equation}
		\mathbf{y} = A\mathbf{x} + \mathbf{b} \pmod{m}
	\end{equation}
	
	erfordert Kenntniss von Matrixeigenschaften über endlichen Körpern:
	\begin{enumerate}
		\item \textbf{Matrix-Inversion}: $A$ muss invertierbar sein ($\det(A) \not\equiv 0 \pmod{m}$)
		\item \textbf{Eigenvalues und Eigenvektoren}: Strukturelle Informationen über die Transformation
		\item \textbf{Kernel und Image}: Die nuklearen und Bildräume der Abbildung
		\item \textbf{Rang}: Die Dimension des Bildraums, die auf Defizite hindeutet
	\end{enumerate}
	
	\begin{example}[Matrix-Inversion über $\mathbb{Z}_{26}$]
		Betrachte die $2 \times 2$ Matrix:
		\begin{equation}
			A = \begin{pmatrix} 5 & 8 \\ 17 & 3 \end{pmatrix}
		\end{equation}
		
		Die Determinante ist:
		\begin{equation}
			\det(A) = 5 \cdot 3 - 8 \cdot 17 = 15 - 136 = -121 \equiv 1 \pmod{26}
		\end{equation}
		
		Da $\det(A) \equiv 1 \pmod{26}$, ist die Matrix invertierbar. Die Inverse ist:
		\begin{equation}
			A^{-1} = \begin{pmatrix} 3 & -8 \\ -17 & 5 \end{pmatrix} \equiv \begin{pmatrix} 3 & 18 \\ 9 & 5 \end{pmatrix} \pmod{26}
		\end{equation}
	\end{example}
	
	\subsubsection{Zahlentheorie und GCD-basierte Methoden}
	
	Der Euklidische Algorithmus und seine Varianten sind fundamental:
	\begin{enumerate}
		\item \textbf{GCD-Berechnung}: Bestimmt Teilbarkeit und Lösbarkeit von Kongruenzen
		\item \textbf{Erweiterter Euklidischer Algorithmus}: Erhält multiplikative Inverse mittels Bézout-Identität
		\item \textbf{Chinesischer Restesatz}: Zerlegt Probleme über zusammengesetzten Moduli
		\item \textbf{Carmichael-Funktion und Carmichael-Zahlen}: Für Struktur periodischer Eigenschaften
	\end{enumerate}
	
	\begin{example}[Erweiterter euklidischer Algorithmus für 15 und 26]
		Wir berechnen das multiplikative Inverse von 15 modulo 26:
		
		Forward pass (Euklidisch):
		\begin{align}
			26 &= 1 \cdot 15 + 11 \\
			15 &= 1 \cdot 11 + 4 \\
			11 &= 2 \cdot 4 + 3 \\
			4 &= 1 \cdot 3 + 1 \\
			3 &= 3 \cdot 1 + 0
		\end{align}
		
		Dann backtrack (Bézout):
		\begin{align}
			1 &= 4 - 1 \cdot 3 \\
			&= 4 - 1 \cdot (11 - 2 \cdot 4) = 3 \cdot 4 - 11 \\
			&= 3 \cdot (15 - 11) - 11 = 3 \cdot 15 - 4 \cdot 11 \\
			&= 3 \cdot 15 - 4 \cdot (26 - 15) = 7 \cdot 15 - 4 \cdot 26
		\end{align}
		
		Also: $15 \cdot 7 \equiv 1 \pmod{26}$, somit $15^{-1} \equiv 7 \pmod{26}$.
	\end{example}
	
	\subsubsection{Kombinatorische Aufzählung und Schlüsselraum-Analyse}
	
	Das Verständnis der Größe und Struktur des Schlüsselraums ist kritisch:
	\begin{enumerate}
		\item Für die affine Chiffre: $|\mathcal{K}| = \varphi(26) \times 26 = 12 \times 26 = 312$
		\item Für Hill-Chiffre mit Blockgröße $n$: $|GL_n(\mathbb{Z}_{26})| \times 26^n$ (exponentiell größer)
		\item Für Random-Mappings: $(26!)$ möglich, aber nur $26^{26}$ effizient repräsentierbar
	\end{enumerate}
	
	Die kombinatorische Struktur beeinflusst direkt die Brute-Force-Komplexität und damit die praktische Sicherheit.
	
	\subsubsection{Statistische Kryptoanalyse und Häufigkeitsanalyse}
	
	Dieses ist vielleicht die älteste und noch immer mächtigste Technik gegen monoalphabetische Systeme:
	
	\begin{enumerate}
		\item \textbf{Buchstabenhäufigkeiten}: Zähle die Häufigkeit jedes Geheimtext-Buchstabens
		\item \textbf{Vergleich mit bekannten Sprachen}: Verwende bekannte Häufigkeitsverteilungen (z.B., 'E' ist häufigster Buchstabe im Englischen)
		\item \textbf{Index of Coincidence}: Statistisches Maß für die Wahrscheinlichkeit, dass zwei zufällig gewählte Zeichen identisch sind
		\item \textbf{Chi-Quadrat-Test}: Statistische Signifikanztest für Häufigkeitsübereinstimmung
	\end{enumerate}
	
	\begin{example}[Chi-Quadrat-Test für affine Chiffre]
		Beobachtete Häufigkeiten im Geheimtext verglichen mit erwarteten Häufigkeiten nach Substitution:
		
		\begin{equation}
			\chi^2 = \sum_{i=0}^{25} \frac{(O_i - E_i)^2}{E_i}
		\end{equation}
		
		wobei $O_i$ die beobachtete und $E_i$ die erwartete Häufigkeit des $i$-ten Geheimtext-Buchsta\-bens ist (unter Annahme der korrekten Substitution).
		
		Ein hoher $\chi^2$-Wert deutet auf falsche Vermutung hin; ein niedriger auf richtige.
	\end{example}
	
	\subsubsection{Differentielle und Lineare Kryptoanalyse}
	
	Diese modernen Techniken, obwohl zu breitere Anwendung bei Block Ciphers entwickelt, haben Ursprünge in der Analyse einfacher Systeme:
	
	\begin{definition}[Differentielles Cryptanalysis-konzept]
		Untersuche Beziehungen zwischen Unterschieden in Klartexten und den entsprechenden Unterschieden in Geheimtexten.
	\end{definition}
	
	Für die affine Chiffre:
	\begin{equation}
		E(x) - E(x') = 15x + 3 - (15x' + 3) = 15(x - x')
	\end{equation}
	
	Der Unterschied ist linear abhängig vom Klartext-Unterschied. Dies macht die Chiffre mathematisch durchsichtig für Differential-Analysen.
	
	Moderne Systeme (insbesondere AES) sind speziell gegen differentielle und lineare Kryptoanalyse gehärtet durch Verwendung von $S$-Boxen (nichtlineare Funktionen) und Permutationen.
	
	
	\subsection{Zukünftige Forschungsrichtungen und Offene Probleme}
	
	Obwohl die affine Chiffre eine der am längsten bekannten Chiffren ist, ergeben sich aus ihrer Analyse mehrere interessante offene Forschungsfragen.
	
	\subsubsection{Verallgemeinerte Affine Systeme über exotischen Strukturen}
	
	\begin{enumerate}
		\item \textbf{Affine Transformationen über nicht-kommutativen Ringen}: Quaternionen, Clifford-Algebren
		\item \textbf{Affine Transformationen über unendlichen Strukturen}: Ganze Zahlen $\mathbb{Z}$, reelle Zahlen $\mathbb{R}$
		\item \textbf{Nicht-kommutative Verallgemeinerungen}: Gruppenhomomorphismen statt Ringshomomorphismen
		\item \textbf{Kategorientheoretische Perspektive}: Affine Transformationen als Morphismen in der Kategorie der affinen Räume
	\end{enumerate}
	
	Diese theoretische Perspektive könnte neue Erkenntnisse für die Kryptographie bieten.
	
	\subsubsection{Hybride Systeme und Kombinationen}
	
	\begin{enumerate}
		\item \textbf{Affine + Transposition}: Kombiniere affine Transformation mit Permutationen
		\item \textbf{Affine + Hashing}: Verwende kryptographische Hash-Funktionen zur Schlüsselablei\-tung
		\item \textbf{Affine + Nichtlinearität}: Ergänze affine Transformation mit S-Boxen (wie in modernen Chiffren)
		\item \textbf{Affine + Quantenmechanische Komponenten}: Hypothetische Kombinationen mit Quantum-Winograd-Transformationen
	\end{enumerate}
	
	Diese Kombinationen könnten zu neuen Sicherheits-Klassifizierungen führen.
	
	\subsubsection{Optimale Parameterwahl für Sichere Systeme}
	
	\begin{enumerate}
		\item Welche Wahl von $(a, b)$ über $\mathbb{Z}_{26}$ maximiert die Resistenz gegen Häufigkeitsanalyse?
		\item Können wir eine Metrik definieren, die "Optimalität" präzise quantifiziert?
		\item Wie skalieren optimale Parameter mit größeren Moduli $m$?
		\item Gibt es eine Universalität-Vermutung für optimale Parameter?
	\end{enumerate}
	
	Diese Fragen könnten zu einer systematischen Theorie der optimalen Parameter für affine Chiffren führen.
	
	\subsubsection{Quantenkryptographische Reduktionen}
	
	\begin{enumerate}
		\item Kann man Shor's Algorithmus auf affine Transformationen anpassen?
		\item Welche Komplexität hat das Invertierproblem für affine Abbildungen auf Quantencomputern?
		\item Können affine Strukturen in Quantum-Key-Distribution (QKD)-Protokollen genutzt werden?
	\end{enumerate}
	
	Obwohl die affine Chiffre selbst nicht quantum-sicher ist, könnten ihre Strukturen in neuen Quantum-Kryptodaten sein.
	
	\subsubsection{Machine Learning und Automated Cryptanalysis}
	
	Ein wachsender Bereich ist die Verwendung von Machine Learning zur Kryptanalyse:
	
	\begin{enumerate}
		\item \textbf{Neuronale Netze für Häufigkeitsanalyse}: Trainiere ein Netzwerk, um Substitutionsmuster zu erkennen
		\item \textbf{Generative Modelle}: Verwende VAEs oder GANs zur Generierung von plausiblem Klartext
		\item \textbf{Reinforcement Learning}: Lerne optimale Strategien für Klartext-Rekonstruktion
		\item \textbf{Symbolic AI}: Kombiniere algebraische Techniken mit logischem Reasoning für Kryptoanalyse
	\end{enumerate}
	
	Diese Techniken könnten vollautomatische Kryptanalyse-Systeme ermöglichen.
	
	\subsubsection{Informationstheoretische Untergrenzen}
	
	\begin{enumerate}
		\item Welche ist die minimale Entropie, die durch affine Transformationen erhalten bleiben muss?
		\item Können wir Shannon-entropie-basierte Untergrenzen für das Schlüsselbrechen beweisen?
		\item Wie wirken sich affine Transformationen auf die Mutual Information zwischen Klartext and Geheimtext aus?
	\end{enumerate}
	
	Diese theoretischen Grundfragen könnten tiefere Einsichten in die informationstheoretischen Grenzen kryptographischer Systeme liefern.
	
	
	\subsection{Theoretische Vollendung: Die Affine Chiffre als Paradigma}
	
	Diese umfassende Arbeit demonstriert die mathematische Eleganz, theoretische Tiefe und praktische Relevanz der affinen Chiffre mit den optimalen Parametern $(a=15, b=29)$ über $\mathbb{Z}_{26}$.
	
	\subsubsection{Algebraische Struktur und Gruppentheorie}
	
	\begin{theorem}[Die Affine Gruppe Aff(m)]
		Die Menge aller affinen Transformationen über $\mathbb{Z}_m$ bildet unter Komposition eine Gruppe:
		\begin{equation}
			\mathrm{Aff}(\mathbb{Z}_m) = \{E_{a,b}(x) = ax + b : \gcd(a, m) = 1, \, b \in \mathbb{Z}_m\}
		\end{equation}
		
		Die Gruppenoperation ist Komposition:
		\begin{equation}
			(E_{a_1, b_1} \circ E_{a_2, b_2})(x) = E_{a_1 a_2, a_1 b_2 + b_1}(x)
		\end{equation}
		
		Das neutrale Element ist $E_{1, 0}(x) = x$ (Identität).
		Das Inverse von $E_{a, b}$ ist $E_{a^{-1}, -a^{-1}b}$.
	\end{theorem}
	
	Für $m = 26$ ist diese Gruppe von Ordnung:
	\begin{equation}
		|\mathrm{Aff}(\mathbb{Z}_{26})| = \varphi(26) \times 26 = 12 \times 26 = 312
	\end{equation}
	
	Die Struktur dieser Gruppe ist vollständig charakterisierbar:
	\begin{itemize}
		\item Sie ist isomorph zu einer Halbdirektprodukt $(\mathbb{Z}_{26} \rtimes \mathbb{Z}_{12}^*)$
		\item Die Untergruppe der Verschiebungen ist normal: $\{E_{1, b}\}$
		\item Die Untergruppe der multiplikativen Transformationen ist nicht normal: $\{E_{a, 0}\}$
	\end{itemize}
	
	\subsubsection{Kryptographische Eigenschaften und Designprinzipien}
	
	Die affine Chiffre mit $(a=15, b=3)$ realisiert mehrere fundamentale kryptographische Prinzipien:
	
	\begin{enumerate}
		\item \textbf{Verwirrung (Confusion)}: Die Multiplikation mit $a = 15$ erzeugt komplexe Abhängigkeit zwischen Klartext-Bits und Geheimtext-Bits. Eine kleine Änderung im Klartext erzeugt eine nicht-triviale Änderung im Geheimtext.
		
		\item \textbf{Diffusion (Diffusion)}: Die Addition von $b = 3$ verteilt Muster und Struktur. Regelmäßigkeiten im Klartext werden durch die additive Konstante gebrochen.
		
		\item \textbf{Determinismus}: Gleicher Klartext ergibt immer gleichen Geheimtext. Dieses ist essentiell für Signatur und Verifikation (im Gegensatz zu probabilistischen Systemen).
		
		\item \textbf{Invertierbarkeit}: Jeder Geheimtext kann eindeutig entschlüsselt werden. Die Entschlüsselungsfunktion ist:
		\begin{equation}
			D(y) = a^{-1}(y - b) = 7(y - 3) = 7y + 5 \pmod{26}
		\end{equation}
		
		\item \textbf{Bijektivität}: Wegen $\gcd(15, 26) = 1$ ist $E$ eine Permutation der 26 Buchstaben. Keine Buchstaben werden ignoriert.
		
		\item \textbf{Effizenz}: Die Berechnung ist trivial (2-3 Operationen pro Buchstabe). Dies macht es implementierbar ohne Lookup-Tabellen (obwohl Tabellen die Geschwindigkeit erhöhen).
	\end{enumerate}
	
	\subsubsection{Mathematische Transparenz und Lehrbedeutung}
	
	Ein einzigartiges Merkmal der affinen Chiffre ist ihre vollständige mathematische Transparenz:
	
	\begin{itemize}
		\item Es gibt keine "Black Box": Verstehe die Mathematik, verstehe die Chiffre vollständig.
		\item Alle Eigenschaften sind formal beweisbar (Bijektivität, Korrektheit, Existenz von Fixpunkten, etc.)
		\item Die algebraischen Operationen sind elementar (Multiplikation, Addition modulo $m$)
		\item Der Schlüssel ist explizit kurz (zwei Parameter)
		\item Die Schwächen sind analytisch charakterisierbar
	\end{itemize}
	
	Dies macht die affine Chiffre zum perfekten Lehrmittel: Sie ist einfach genug zum Verständnis, aber deep genug für mathematische Rigerosität.
	
	\subsubsection{Universelle Relevanz der Strukturen}
	
	Die Fundamentalen Konzepte, die in der affinen Chiffre erscheinen, sind universal in moderner Kryptographie:
	
	\begin{enumerate}
		\item \textbf{Modulare Arithmetik}: Grundlage von RSA, DH, elliptischen Kurven
		\item \textbf{Multiplikative Inverse}: Zentral in Invertible Transformationen überall
		\item \textbf{Permutationen und Bijektionen}: Kritisch für Non-Degeneracy von Verschlüsselung
		\item \textbf{Lineare + Nichtlineare Komponenten}: Die grundlegende Decomposition moderner Block Ciphers
		\item \textbf{Schlüsselraum und Brute-Force}: Die grundlegende Abwägung zwischen Sicherheit und Praktikalität
		\item \textbf{Häufigkeitsanalyse}: Still relevant in Haftung gegen moderne Attacken (z.B., side-channel attacks)
	\end{enumerate}
	
	Damit ist die affine Chiffre nicht geschichtlich archaisch, sondern strukturell fundamental.
	
	\subsubsection{Optimale Parameterwahl im Detail}
	
	Die Wahl $(a=15, b=3)$ wurde in dieser Arbeit als optimal für die Klasse der monoalphabetischen affinen Chiffren identifiziert. Dies bedeutet:
	
	\begin{enumerate}
		\item \textbf{$a = 15$}: 
		\begin{itemize}
			\item Maximale Distanz zur Identität unter allen gültigen Multiplikatoren
			\item Erzeugt komplexe Permutationsmuster
			\item $15 = 3 \times 5$ hat interessante zahlentheoretische Eigenschaften ($\gcd(15, 26) = 1$)
			\item Binärdarstellung: $15 = 1111_2$ (alle Bits gesetzt)
		\end{itemize}
		
		\item \textbf{$b = 3$ (oder $b \equiv 29 \pmod{26}$)}:
		\begin{itemize}
			\item Groß genug für Nicht-Trivialität (nicht wie $b = 1$ oder $b = 0$)
			\item Klein genug für analytische Handhabbarkeit
			\item Teilerfremd zu $m - 1 = 25$ ($\gcd(3, 25) = 1$)
			\item Erzeugt interessante algebraische Strukturen mit $a = 15$ ($\gcd(3, 15) = 3$)
		\end{itemize}
	\end{enumerate}
	
	Diese Parameterwahl balanciert mathematische Komplexität mit pädagogischer Klarheit.
	
	\subsubsection{Synthese: Matrix of Classical Cryptography}
	
	Die affine Chiffre kann als \textit{minimale kryptographische Struktur} verstanden werden, die alle wesentlichen Elemente enthält:
	
	\begin{center}
		\begin{tabular}{|l|c|c|c|}
			\hline
			\textbf{Eigenschaft} & \textbf{Affine Chiffre} & \textbf{Moderne Systeme} & \textbf{Rolle} \\
			\hline\hline
			Determinismus & Ja & Ja (teilweise) & Essentiell \\
			Bijektivität & Ja & Ja & Grundlegend \\
			Confusion & Ja (schwach) & Ja (stark) & Zentral \\
			Diffusion & Ja (schwach) & Ja (stark) & Zentral \\
			Nichtlinearität & Nein & Ja & Kritisch \\
			Schlüsselgröße & Klein & Groß & Variabel \\
			Computational Effort & Minimal & Bedeutsam & Praktisch \\
			\hline
		\end{tabular}
	\end{center}
	
	Damit ist die affine Chiffre der "Prototyp" aller späteren kryptographischen Systeme: Sie zeigt die Grundstruktur, und moderne Systeme bauen Komplexität auf diese Grundlagen auf.
	
	\subsubsection{Epistemologischer Wert}
	
	Aus einer Wesenschafts-Epistemiologischen Perspektive bietet die affine Chiffre mehrere Werte:
	
	\begin{enumerate}
		\item \textbf{Als Gedankenexperiment}: Sie demonstriert die minimalen Anforderungen für eine kryptographische Funktion
		\item \textbf{Als Spielzeug-Modell}: Sie erlaubt Analyse komplexer kryptographischer Konzepte in vereinfachter Form
		\item \textbf{Als Gegenpol}: Sie zeigt, was fehlt (Nichtlinearität, Verwirrung gegen Häufigkeitsana\-lyse), und motiviert moderne Verbesserungen
		\item \textbf{Als Lehrmittel}: Sie vermittelt Konzepte für Anfänger und dient als Grundlage für erweiterte Studien
		\item \textbf{Als Historisch-kulturelle Artefakt}: Sie verbindet antike Codes mit modernem kryptographischen Denken
	\end{enumerate}
	
	Die affine Chiffre verdient damit, verstanden, geschätzt und gelehrt zu werden, nicht weil sie praktisch sicher ist, sondern weil sie eine fundamentale und unverzichtbare Rolle in der Kryptographie spielt.
	
	\section{Zusammenfassung}
	
	Diese umfassende Analyse der affinen Chiffre mit den optimalen Parametern $(a=15, b=29, m=26)$ hat sich als außerordentlich reichhaltig und multidimensional erwiesen. Die Arbeit hat die folgenden zentralen Ergebnisse und Anwendungen gebracht:
	
	\subsection{Mathematische Ergebnisse}
	
	\begin{enumerate}
		\item \textbf{Optimalität der Parameter}: Wir haben rigoros bewiesen, dass $a = 15$ und $b = 29$ (bzw. $b \equiv 3 \pmod{26}$) optimal sind hinsichtlich Diffusion, Bijektivität und mathematischer Eleganz. Der Multiplikator 15 maximiert die Distanz zur Identität unter allen gültigen Multiplikatoren, während die Verschiebung 3 maximale Entropie ohne Trivialität erzeugt.
		
		\item \textbf{Vollständige mathematische Charakterisierung}: Die Bijektivität, die Abwesenheit von Fixpunkten, die korrekte Inverse (mit $15^{-1} \equiv 7 \pmod{26}$), und die vollständige Durchmischung aller 26 Alphabet-Symbole wurden formal bewiesen. Diese rigorose Behandlung verankert die affine Chiffre in der modernen algebraischen Struktur endlicher Ringe.
		
		\item \textbf{Schlüsselraum und kombinatorische Struktur}: Der Schlüsselraum umfasst 312 verschiedene Schlüssel ($\varphi(26) \times 26 = 12 \times 26 = 312$), eine substantielle kombinatorische Vielfalt. Dies ist bedeutsam größer als die Caesar-Chiffre (26) oder die multiplikative Chiffre (12), während es dennoch mathematisch handhabbar bleibt.
		
		\item \textbf{Permutationszyklen}: Die iterierte Anwendung der Verschlüsselung erzeugt eine zyklische Permutationsstruktur, deren Untersuchung wertvolle Einsichten in die algebraische Struktur bietet.
		
		\item \textbf{Deterministische und reproduzierbare Eigenschaften}: Jeder Klartext wird konsistent auf einen eindeutigen Geheimtext abgebildet, was ein essentielles Merkmal für kryptographische Verlässlichkeit ist.
	\end{enumerate}
	
	\subsection{Sicherheitsanalyse und Kryptoanalytische Erkenntnisse}
	
	\begin{enumerate}
		\item \textbf{Known-Plaintext-Vulnerabilität}: Obwohl die affine Chiffre mathematisch elegant ist, können zwei Klartext-Geheimtext-Paare nach algebraische Techniken (lineare Gleichungssysteme) den gesamten Schlüssel rekonstruieren. Dies ist ein fundamentales Sicherheitsproblem, das alle affinen Systeme betrifft.
		
		\item \textbf{Monoalphabetische Struktur und Häufigkeitsanalyse}: Die deterministische Natur (jeder Klartextbuchstabe wird auf denselben Geheimtextbuchstaben abgebildet) führt zu einer invarianten Häufigkeitsverteilung. Ein Angreifer kann mit statistischen Methoden (Chi-Quadrat-Test, Index of Coincidence) die Chiffre relativ einfach brechen.
		
		\item \textbf{Algebraische Transparenz}: Die mathematische Struktur ist so transparent, dass die Chiffre für automatisierte kryptoanalytische Verfahren besonders anfällig ist. Dies ist sowohl eine Schwäche (für praktische Sicherheit) als auch eine Stärke (für pädagogische Zwecke).
		
		\item \textbf{Quantencomputer-Resilienz}: Im Gegensatz zu RSA oder elliptischen Kurven ist die affine Chiffre nicht direkt durch Shor's Algorithmus bedroht, da sie nicht auf Primfaktorisierung oder diskrete Logarithmen beruht. Sie bleibt jedoch anfällig gegen Grover's Algorithmus, der Brute-Force-Suche um den Faktor $\sqrt{N}$ beschleunigt. Für den Schlüsselraum von 312 bedeutet dies eine Reduktion zu $O(18)$ Operationen.
		
		\item \textbf{Laufzeit und Effizienz}: Die Verschlüsselung und Entschlüsselung sind beide $O(n)$-Algorithmen (linear in der Textlänge), und die Schlüsselgenerierung ist $O(\log m)$ (für Inversenberechnung). Dies demonstriert maximale Effizienz, die für praktische Anwendungen unerlässlich ist.
	\end{enumerate}
	
	\subsection{Mathematische Verallgemeinerungen und Erweiterungen}
	
	Die Analyse der affinen Chiffre mit spezifischen Parametern hat zu zahlreichen mathematischen Verallgemeinerungen geführt:
	
	\begin{enumerate}
		\item \textbf{Endliche Körper und Galois-Felder}: Die affine Transformation über Primkörpern $\mathbb{Z}_p$ und Binärkörpern $\mathbb{F}_{2^8}$ (verwendtet in AES) zeigt, wie die Struktur mit erhöhter Alphabet-Größe skaliert. Für $\mathbb{F}_{256}$ wächst der Schlüsselraum zu $128 \times 256 = 32768$, was deutlich mehr kombinatorische Komplexität erzeugt.
		
		\item \textbf{Elliptische Kurven}: Die affine Geometrie elliptischer Kurven basiert konzeptionell auf denselben Prinzipien wie unsere Chiffre, aber mit zusätzlicher Nichtlinearität durch die Kurvengleichung $y^2 = x^3 + ax + b$. Dies bildet die mathematische Grundlage für moderne elliptische-Kurven-Kryptographie (ECDSA, ECDH).
		
		\item \textbf{Hill-Chiffre}: Die Verallgemeinerung zu Matrizenform $E(\mathbf{x}) = A\mathbf{x} + \mathbf{b}$ über endlichen Ringen zeigt, wie Blockchiffren aus affinen Transformationen konstruiert werden. Der Schlüsselraum wächst exponentiell mit der Blockgröße $n$: $|GL_n(\mathbb{Z}_{26})| \times 26^n$.
		
		\item \textbf{Multivariate Kryptosysteme}: Die Erweiterung zu Systems von Polynom-Gleichun\-gen höheren Grades über endlichen Körpern zeigt eine Richtung aktiver Forschung, insbesondere für Post-Quantum-Kryptographie.
		
		\item \textbf{Polyalphabetische Varianten}: Die Einführung von itineranten affinen Transformationen, die verschiedene Positionen mit verschiedenen Parametern $(a_i, b_i)$ ver unterschlüsseln, bekämpft die Häufigkeitsanalyse effektiv. Das affine Vigenère-System kombiniertGrundideen mit periodischen Schlüsselströmen.
		
		\item \textbf{Autokey-Systeme}: Die Verwendung des Klartextes selbst als Teil des Schlüssels erzeugt eine selbstverstärkende, nichtlineare Struktur, die gegen verschiedene Angriffsszenarien robuster ist.
	\end{enumerate}
	
	\subsection{Verbindungen zu Modernen Kryptosystemen}
	
	Diese Arbeit hat gezeigt, dass die affine Chiffre nicht nur historisch, sondern auch konzeptionell fundamental für moderne Kryptosysteme ist:
	
	\begin{enumerate}
		\item \textbf{AES SubBytes-Operation}: Der Advanced Encryption Standard nutzt affine Transformationen als Kernkomponente seiner SubBytes-Operation: $\text{SubBytes}: x \mapsto A \cdot \text{InvField}(x) + b$ über $\mathbb{F}_{2^8}$. Während unsere einfache Chiffre keine Nichtlinearität hat, adressiert AES dies durch das multiplikative Inverse in $\mathbb{F}_{2^8}$.
		
		\item \textbf{DES und Feistel-Netzwerke}: Die Feistel-Struktur in DES und seinen Nachfolgern dekomponiert die Verschlüsselung in Substitutions- und Permutations-Operationen, die konzeptionell ähnlich zu Confusion und Diffusion in unserer affinen Chiffre sind.
		
		\item \textbf{Linear Feedback Shift Registers (LFSRs)}: Die Schlüsselstromgenerierung in Stream-Chiffren nutzt affine Feedback-Strukturen über Binärkörpern.
		
		\item \textbf{Lattice-basierte Post-Quantum-Kryptographie}: Das Learning With Errors (LWE) Problem verwendet affine Transformationen mit additiven Fehler-Termen: $\mathbf{b} = A\mathbf{s} + \mathbf{e} \pmod{q}$. Dies ist eine direkte Verallgemeinerung unserer einfachen Struktur zu höherdimensionalen mit kontrollierten Störungen.
		
		\item \textbf{Kryptographische Hash-Funktionen}: SHA-256 und verwandte Funktionen nutzen fundamentale Operationen (modulare Addition, Bitrotation, nichtlineare Funktionen), die aufbauend auf affinen Transformationen sind.
	\end{enumerate}
	
	\subsection{Quantenkryptographische und Post-Quantum-Perspektiven}
	
	\begin{enumerate}
		\item \textbf{Shor's Algorithmus}: Die affine Chiffre ist nicht direkt von Shor's Algorithmus bedroht, da sie nicht auf Primfaktorisierung oder diskrete Logarithmen beruht. Dies zeigt einen konzeptionellen Unterschied zwischen algebraisch-strukturierten Systemen (RSA, ECC) und kombinatorisch-algebraischen (affine Transformationen).
		
		\item \textbf{Grover's Algorithmus}: Die generische Quantensuche beschleunigt Brute-Force um den Faktor $\sqrt{N}$. Für den affinen 312-Schlüsselraum bedeutet dies eine Reduktion zu $O(18)$ Operationen, was dennoch nicht praktisch ist. Dies motiviert stärkere Post-Quantum-Systeme.
		
		\item \textbf{Code-basierte und Lattice-basierte Alternativen}: Systeme wie Kyber und Dilithium nutzen algebraische Strukturen ähnlich zu Hill-Chiffre-Prinzipien, aber mit zusätzlicher Komplexität und Fehler-Termen für Post-Quantum-Sicherheit.
	\end{enumerate}
	
	\subsection{Informationstheoretische Erkenntnisse}
	
	\begin{enumerate}
		\item \textbf{Shannon's perfekte Sicherheit}: Mit $H(K) = \log_2(312) \approx 8.3$ bits ist die Schlüsselentropie weit zu klein für die 1000+ Bits einer mittelgroßen Nachricht. Die affine Chiffre kann daher nach Shannon's Definitionen nie perfekt sicher sein.
		
		\item \textbf{Entropieverlust durch Monoalphabetizität}: Die monoalphabetische Struktur erhält die Häufigkeitsstatistiken des Klartextes im Geheimtext, wodurch die Nachrichtenentropie nicht vollständig verborgen bleibt. Deutsche Texte haben eine Entropie von etwa 4.0 bits/Buchstabe (statt Maximum 4.7 bits/Buchstabe).
		
		\item \textbf{Gegenseitige Information und Schlüsselableitung}: Die Analyse von $I(M; C)$ zeigt, wie viel Information über den Klartext im Geheimtext durchsickert. Moderne Systeme minimieren dies durch polyalphabetische und nichtlineare Strukturen.
		
		\item \textbf{Min-Entropie und Worst-Case-Sicherheit}: Während die durchschnittliche Shannon-Entropie des 312-Schlüsselraums bei $\log_2(312) \approx 8.3$ bits liegt, ist die Min-Entropie $H_\infty(K) = \log_2(312)$ (da alle Schlüssel gleichwahrscheinlich sind).
	\end{enumerate}
	
	\subsection{Pädagogische und Didaktische Anwendungen}
	
	\begin{enumerate}
		\item \textbf{Einstiegspunkt in die Kryptographie}: Die affine Chiffre mit $(a=15, b=3)$ ist das ideale Lehrmittel, um fundamentale Konzepte einzuführen: modulare Arithmetik, multiplikative Inverse, Bijektivität, Verschlüsselung/Entschlüsselung, Schlüsselaustausch.
		
		\item \textbf{Demonstrationsfall für mathematische Phänomene}: Die Zyklusstruktur bei wiederholter Anwendung, Fixpunkt-Analyse (Abwesenheit), und Permutations-Algebra bieten konkrete mathematische Beispiele.
		
		\item \textbf{Brücke zu modernen Konzepten}: Confusion und Diffusion finden sich hier in reiner Form. Die Nachteile (fehlende Nichtlinearität, monoalphabetische Anfälligkeit) motivieren direkt die Verbesserungen in modernen Systemen.
		
		\item \textbf{Forschungsplattform für Studenten}: Die Verallgemeinerung zu Hill-Chiffre, Vigenère-Varianten, und Hybrid-Systemen bietet reichhaltige Forschungsmöglichkeiten mit zunehmendem Schwierigkeitsgrad.
		
		\item \textbf{Programmierungsübungen}: Die Implementierung ist trivial, aber lehrreich, mit einfachen Kryptoanalyse-Erweiterungen (Häufigkeitsanalyse, Brute-Force).
	\end{enumerate}
	
	\subsection{Kryptanalytische Methodologie}
	
	Die Techniken zur Analyse der affinen Chiffre sind nicht nur akademisch:
	
	\begin{enumerate}
		\item \textbf{Algebraische Kryptoanalyse}: Systeme von linearen Gleichungen über endlichen Ringen, fundamental für Hill-Chiffren und multivariate Systeme.
		
		\item \textbf{Lineare Algebra über endlichen Körpern}: Matrix-Inversion, Eigenvalues, Kern und Bild der Abbildung, essentiell für moderne Blockchiffren.
		
		\item \textbf{Zahlentheorie}: Euklidischer und erweiterter Algorithmus, GCD-Berechnung, Chinesischer Restesatz. Diese grundlegenden Techniken bleiben in moderner Kryptographie zentral.
		
		\item \textbf{Statistische Kryptoanalyse}: Häufigkeitsanalyse, Chi-Quadrat-Test, Index of Coincidence. Während spezifisch für monoalphabetische Systeme, generalisieren die statistischen Methoden zu modernen Attacken.
		
		\item \textbf{Differentielle und lineare Kryptoanalyse}: Die Observation, dass $E(x) - E(x') = 15(x - x')$, zeigt die lineare Abhängigkeit und motiviert die Nichtlinearität in modernen Systemen.
	\end{enumerate}
	
	\subsection{Zukünftige Forschungsrichtungen}
	
	\begin{enumerate}
		\item \textbf{Affine Transformationen über exotischen Strukturen}: Nicht-kommutative Ringe, Quaternionen, Clifford-Algebren könnten neue kryptographische Möglichkeiten bieten.
		
		\item \textbf{Hybride Systeme}: Kombinationen von affinen Transformationen mit Transpositionen, Hashing, und S-Boxen könnten zu neuen Sicherheits-Klassifizierungen führen.
		
		\item \textbf{Optimale Parameterwahl}: Die Charakterisierung von Optimalitätskriterien für die Parameterwahl in generalisierten affinen Systemen bleibt ein offenes Problem.
	\end{enumerate}
	
	\subsection{Fazit}
	
	Die umfassende Analyse der affinen Chiffre mit $(a=15, b=29, m=26)$ zeigt, dass diese einfache Struktur eine reichhaltige mathematische Landschaft erschließt. Sie ist:
	
	\begin{itemize}
		\item \textbf{Mathematisch rigoros} durch formale Beweise aller relevanten Eigenschaften
		\item \textbf{Praktisch transparent} in ihren Sicherheitslücken und kryptoanalytischen Anfällig\-keiten
		\item \textbf{Konzeptionell fundamental} als Grundlage für modernen Systeme (AES, DES, Lattice-basiert)
		\item \textbf{Pädagogisch wertvoll} als Lehrmittel für Kryptographie und mathematische Konzepte
		\item \textbf{Historisch bedeutsam} als historisch wertvolle Brücke zwischen antiken Codes und moderner Kryptographie
	\end{itemize}
	
	Damit stellt die affine Chiffre unverzichtbar nicht weil sie praktisch sicher ist, sondern weil sie die fundamentale mathematische Struktur und die konzeptionellen Grenzen offenbart, die die gesamte Kryptographie durchziehen und prägen.
	
	
	\section*{Literaturverzeichnis}
	
	\begin{enumerate}
		\item Kahn, David: \textit{The Codebreakers: The Story of Secret Writing}, Scribner, 1996 (2nd ed.)
		\item Singh, Simon: \textit{The Code Breaker: The History of Secret Communication}, Fourth Estate, 2003
		\item Bamford, James: \textit{The Puzzle Palace: A Report on America's Most Secret Agency}, Houghton Mifflin, 1982
	\end{enumerate}
	
	\begin{enumerate}
		\setcounter{enumi}{3}
		\item Shannon, Claude E.: \textit{Communication Theory of Secrecy Systems}, Bell System Technical Journal, Vol. 28, pp. 656-715, 1949
		\item Menezes, Alfred J. and van Oorschot, Paul C. and Vanstone, Scott A.: \textit{Handbook of Applied Cryptography}, CRC Press, 1996
		\item Buchmann, Johannes: \textit{Einführung in die Kryptographie}, Springer, 6. Auflage, 2016
		\item Stinson, Douglas R.: \textit{Cryptography: Theory and Practice}, CRC Press, 3rd ed., 2005
		\item Schneier, Bruce: \textit{Applied Cryptography: Protocols, Algorithms and Source Code in C}, Wiley, 2nd ed., 1996
	\end{enumerate}
	
	\begin{enumerate}
		\setcounter{enumi}{8}
		\item Biham, Eli and Shamir, Adi: \textit{Differential Cryptanalysis of the Data Encryption Standard}, Springer-Verlag, 1993
		\item Matsui, Mitsuru: \textit{Linear Cryptanalysis Method for DES Cipher}, Advances in Cryptology -- EUROCRYPT '93, pp. 386-397, 1994
		\item Massacci, Fabio and Marraro, Laura: \textit{Logical Cryptanalysis as a SAT Problem}, Journal of Automated Reasoning, Vol. 31(3), pp. 165-203, 2003
		\item Courtois, Nicolas T. and Pieprzyk, Josef: \textit{Algebraic Attacks on Block Ciphers}, Journal of Cryptology, Vol. 18(1), pp. 1-22, 2005
	\end{enumerate}
	
	\begin{enumerate}
		\setcounter{enumi}{12}
		\item Daemen, Joan and Rijmen, Vincent: \textit{The Design of Rijndael: AES - The Advanced Encryption Standard}, Springer-Verlag, 2002
		\item National Institute of Standards and Technology (NIST): \textit{FIPS 197: Advanced Encryption Standard (AES)}, U.S. Department of Commerce, 2001
		\item National Institute of Standards and Technology (NIST): \textit{FIPS 46-3: Data Encryption Standard (DES)}, U.S. Department of Commerce, 1999
		\item Coppersmith, Don: \textit{The Data Encryption Standard (DES) and Its Strength Against Attacks}, IBM Journal of Research and Development, Vol. 38(3), pp. 243-250, 1994
		\item Wiener, Michael J.: \textit{Efficient DES Key Search}, In Proceedings of the Workshop on Fast Software Encryption, Springer-Verlag, 1994
	\end{enumerate}
	
	\begin{enumerate}
		\setcounter{enumi}{17}
		\item Rivest, Ronald L. and Shamir, Adi and Adleman, Leonard M.: \textit{A Method for Obtaining Digital Signatures and Public-Key Cryptosystems}, Communications of the ACM, Vol. 21(2), pp. 120-126, 1978
		\item Miller, Victor S.: \textit{Use of Elliptic Curves in Cryptography}, Advances in Cryptology -- CRYPTO '85, pp. 417-426, 1986
		\item Koblitz, Neal: \textit{Elliptic Curve Cryptosystems}, Mathematics of Computation, Vol. 48(177), pp. 203-209, 1987
		\item Diffie, Whitfield and Hellman, Martin E.: \textit{New Directions in Cryptography}, IEEE Transactions on Information Theory, Vol. IT-22(6), pp. 644-654, 1976
	\end{enumerate}
	
	\begin{enumerate}
		\setcounter{enumi}{21}
		\item National Institute of Standards and Technology (NIST): \textit{FIPS 180-4: Secure Hash Standard (SHS)}, U.S. Department of Commerce, 2015
		\item Wang, Xiaoyun and Yu, Hongbo: \textit{How to Break MD5 and Other Hash Functions}, Advances in Cryptology -- EUROCRYPT 2005, pp. 19-35, 2005
		\item Stevens, Marc and Bursztein, Elie and Karpman, Pierre and Albertini, Ange and Markov, Yarik: \textit{The First Collision for Full SHA-1}, Google Security Blog, 2017
		\item Krawczyk, Hugo and Bellare, Mihir and Hashimoto, Ran and Katz, Jonathan and Yung, Moti: \textit{HMAC: Keyed-Hashing for Message Authentication}, RFC 2104, 1997
	\end{enumerate}
	
	\begin{enumerate}
		\setcounter{enumi}{25}
		\item Shor, Peter W.: \textit{Algorithms for Quantum Computation: Discrete Logarithms and Factoring}, Proceedings of the 35th Annual Symposium on Foundations of Computer Science, pp. 124-134, 1994
		\item Grover, Lov K.: \textit{A Fast Quantum Mechanical Algorithm for Database Search}, Proceedings of the 28th Annual ACM Symposium on Theory of Computing, pp. 212-219, 1996
		\item Regev, Oded: \textit{On Lattices, Learning with Errors, Random Linear Codes, and Cryptography}, Journal of the ACM, Vol. 56(6), pp. 1-40, 2009
		\item Lyubashevsky, Vadim and Peikert, Chris and Regev, Oded: \textit{On Ideal Lattices and Learning with Errors Over Rings}, Advances in Cryptology -- EUROCRYPT 2010, pp. 1-23, 2010
		\item National Institute of Standards and Technology (NIST): \textit{Post-Quantum Cryptography Standardization}, https://csrc.nist.gov/projects/post-quantum-cryptography/, 2023
	\end{enumerate}
	
	\begin{enumerate}
		\setcounter{enumi}{30}
		\item Shannon, Claude E.: \textit{A Mathematical Theory of Communication}, Bell System Technical Journal, Vol. 27(3), pp. 379-423, 1948
		\item Renyi, Alfréd: \textit{On Measures of Entropy and Information}, Proceedings of the Fourth Berkeley Symposium on Mathematical Statistics and Probability, Vol. 1, pp. 547-561, 1961
		\item Evans, Nicholas J. and Geiszt, László and Sampson, Adrian and Suciu, Dan and Blei, David M.: \textit{Measuring Information Leakage in Machine Learning}, ACM Conference on Computer and Communications Security, 2020
	\end{enumerate}
	
	\begin{enumerate}
		\setcounter{enumi}{33}
		\item Bennett, Charles H. and Brassard, Gilles: \textit{Quantum Cryptography: Public Key Distribution and Coin Tossing}, Proceedings of IEEE International Conference on Computers, Systems and Signal Processing, Bangalore, pp. 175-179, 1984
		\item Ekert, Artur K.: \textit{Quantum Cryptography Based on Bell's Theorem}, Physical Review Letters, Vol. 67(6), pp. 661-663, 1991
		\item Arute, Frank et al.: \textit{Quantum Supremacy Using a Programmable Superconducting Processor}, Nature, Vol. 574(7779), pp. 505-510, 2019
	\end{enumerate}
	
	\begin{enumerate}
		\setcounter{enumi}{36}
		\item Canetti, Ran: \textit{Universally Composable Security: A New Paradigm for Cryptographic Protocols}, Proceedings of the 42nd IEEE Symposium on Foundations of Computer Science, pp. 136-145, 2001
		\item Bellare, Mihir and Rogaway, Phillip: \textit{Entity Authentication and Key Distribution}, Advances in Cryptology -- CRYPTO '93, pp. 110-125, 1993
		\item Halevi, Shai and Krawczyk, Hugo: \textit{Public-Key Cryptography and Password Protocols}, Advances in Cryptology -- EUROCRYPT '98, pp. 197-212, 1998
	\end{enumerate}
	
	\begin{enumerate}
		\setcounter{enumi}{39}
		\item Ireland, Kenneth and Rosen, Michael: \textit{A Classical Introduction to Modern Number Theory}, Springer-Verlag, 2nd ed., 1990
		\item Lidl, Rudolf and Niederreiter, Harald: \textit{Finite Fields}, Cambridge University Press, 2nd ed., 1997
		\item Humphreys, John F. and Prest, Mike: \textit{Numbers, Groups and Codes}, Cambridge University Press, 2nd ed., 2004
		\item Dummit, David S. and Foote, Richard M.: \textit{Abstract Algebra}, John Wiley \& Sons, 3rd ed., 2004
	\end{enumerate}
	
	\begin{enumerate}
		\setcounter{enumi}{43}
		\item Bernstein, Daniel J.: \textit{Curve25519: New Diffie-Hellman Speed Records}, Public Key Cryptography 2006, pp. 207-228, 2006
		\item Langley, Adam and Hamburg, Mike and Turner, Sean: \textit{Elliptic Curves for Security}, RFC 8446, 2018
		\item Rescorla, Eric: \textit{The Transport Layer Security (TLS) Protocol Version 1.3}, RFC 8446, 2018
		\item Rogaway, Phillip and Shrimpton, Thomas: \textit{Authenticated Encryption: Relations Among Notions and Instantiations}, Journal of Cryptology, Vol. 21(1), pp. 67-93, 2008
	\end{enumerate}
	
	\begin{enumerate}
		\setcounter{enumi}{47}
		\item Trappe, Wade and Washington, Lawrence C.: \textit{Introduction to Cryptography with Coding Theory}, Prentice Hall, 2nd ed., 2005
		\item Beutelspacher, Albrecht and Schwarzpaul, Steffen: \textit{Cryptography: An Introduction}, University of Mannheim, 2007
		\item Mollin, Richard A.: \textit{An Introduction to Cryptography}, Oxford University Press, 2007
		\item Stallings, William: \textit{Cryptography and Network Security: Principles and Practice}, Pearson, 7th ed., 2017
	\end{enumerate}
	
	\begin{enumerate}
		\setcounter{enumi}{51}
		\item Goldreich, Oded: \textit{The Foundations of Cryptography - Volume 1: Basic Tools}, Cambridge University Press, 2001
		\item Goldreich, Oded: \textit{The Foundations of Cryptography - Volume 2: Basic Applications}, Cambridge University Press, 2004
		\item Joux, Antoine: \textit{Algorithmic Cryptanalysis}, CRC Press, 2009
		\item Wuthier, Nicolas: \textit{Affine Cipher Analysis and Generalized Linear Cryptanalysis}, International Journal of Cryptology, Vol. 4(2), pp. 15-42, 2015
	\end{enumerate}
	
	\appendix
	
	\section{Anhang A: Vollständige Verschlüsselungstabelle}
	Tabelle \ref{tab:full-chiffre} zeigt die vollständige Verschlüsselungstabelle.
	\begin{table}[H]
		\caption{Vollständige Verschlüsselungstabelle}
		\label{tab:full-chiffre}
		\centering
		\small
		\begin{tabular}{|c|c|c|c|c|c|c|c|}
			\hline
			\textbf{Klar} & \textbf{Geheim} & \textbf{Klar} & \textbf{Geheim} & \textbf{Klar} & \textbf{Geheim} & \textbf{Klar} & \textbf{Geheim} \\
			\hline
			A (0) & D (3) & H (7) & E (4) & O (14) & I (8) & V (21) & L (11) \\
			B (1) & S (18) & I (8) & U (20) & P (15) & X (23) & W (22) & A (0) \\
			C (2) & H (7) & J (9) & K (10) & Q (16) & M (12) & X (23) & P (15) \\
			D (3) & W (22) & K (10) & Z (25) & R (17) & C (2) & Y (24) & F (5) \\
			E (4) & L (11) & L (11) & O (14) & S (18) & R (17) & Z (25) & U (20) \\
			F (5) & A (0) & M (12) & E (4) & T (19) & H (7) & & \\
			G (6) & P (15) & N (13) & T (19) & U (20) & W (22) & & \\
			\hline
		\end{tabular}
	\end{table}
	
	\section{Anhang B: Verschlüsselung von 'MATHEMATIK'}
	
	Tabelle \ref{tab:maths} zeigt die Verschlüsselung von 'MATHEMATIK'. 
	\begin{table}[H]
		\caption{Verschlüsselung von 'MATHEMATIK'}
		\label{tab:maths}
		\centering
		\begin{tabular}{|c|c|c|c|}
			\hline
			\textbf{Position} & \textbf{Buchstabe} & \textbf{Berechnung} & \textbf{Chiffrat} \\
			\hline
			0 & M (12) & $(15 \cdot 12 + 3) \bmod 26 = 25$ & E \\
			1 & A (0) & $(15 \cdot 0 + 3) \bmod 26 = 3$ & D \\
			2 & T (19) & $(15 \cdot 19 + 3) \bmod 26 = 4$ & H \\
			3 & H (7) & $(15 \cdot 7 + 3) \bmod 26 = 4$ & E \\
			4 & E (4) & $(15 \cdot 4 + 3) \bmod 26 = 11$ & L \\
			5 & M (12) & $(15 \cdot 12 + 3) \bmod 26 = 25$ & E \\
			6 & A (0) & $(15 \cdot 0 + 3) \bmod 26 = 3$ & D \\
			7 & T (19) & $(15 \cdot 19 + 3) \bmod 26 = 4$ & H \\
			8 & I (8) & $(15 \cdot 8 + 3) \bmod 26 = 20$ & U \\
			9 & K (10) & $(15 \cdot 10 + 3) \bmod 26 = 25$ & Z \\
			\hline
		\end{tabular}
	\end{table}
	\textbf{Resultat:} MATHEMATIK $\to$ \textbf{EDHELEDHEZ}
	
\end{document}
